% =====================================================================
%  EMNLP 2026 -- Paper #3
%  When Does a Chain of Thought Commit?
%  Causal Lookahead Probing of Reasoning-Tuned Language Models
%
%  Master file. Uses the official ACL/EMNLP template (acl.sty) verbatim;
%  template format is NOT modified. Bib: emnlp.bib (acl_natbib.bst).
%  Written sections are filled in.
%  No em dashes anywhere (author preference).
% =====================================================================
\documentclass[11pt]{article}
\usepackage[review]{acl}
\usepackage{times}
\usepackage{latexsym}
\usepackage[T1]{fontenc}
\usepackage[utf8]{inputenc}
\usepackage{microtype}
\usepackage{graphicx}
% ===== Additional packages =====
\usepackage{amsmath,amssymb}
\usepackage{amsthm}
\newtheorem{definition}{Definition}
\newtheorem{theorem}{Theorem}
\newtheorem{proposition}{Proposition}
\newtheorem{lemma}{Lemma}
\newtheorem{corollary}{Corollary}
\newtheorem*{remark}{Remark}
\usepackage{enumitem}
\usepackage{booktabs}
\usepackage{supertabular}
\usepackage{multirow}
\usepackage{xcolor}
\definecolor{ckblue}{HTML}{0072B2}
\definecolor{ckskyblue}{HTML}{56B4E9}
\definecolor{ckgreen}{HTML}{009E73}
\definecolor{ckorange}{HTML}{E69F00}
\definecolor{ckvermillion}{HTML}{D55E00}
\definecolor{ckgray}{HTML}{999999}
\usepackage{algorithm}
\usepackage{algpseudocode}
\usepackage{hyperref}
\usepackage{cleveref}

% ===== Custom commands (shared macros) =====
\newcommand{\CD}{\mathrm{CD}}
\newcommand{\CDP}{\mathrm{CD}_{P}}
\newcommand{\CDM}{\mathrm{CD}_{M}}
\newcommand{\stepk}{k}
\newcommand{\CLP}{\textsc{CLP}}
\newcommand{\CRASPCoT}{\textsc{C-RASP}^{\textsc{CoT}}}
\newcommand{\Pythia}{\textsc{Pythia}}
\newcommand{\Llama}{\textsc{Llama-3.2}}
\newcommand{\Mamba}{\textsc{Mamba}}
\newcommand{\Rdistill}{\textsc{R1-Distill}}

\title{When Does a Chain of Thought Commit?\\
       Causal Lookahead Probing of Reasoning-Tuned Language Models}

\author{Anonymous}

\begin{document}
\maketitle

% =====================================================================
%  ABSTRACT
% =====================================================================
\begin{abstract}
Reasoning models generate long chains of thought before answering, and users
increasingly read those chains as evidence of how a decision was reached. But
when does the model actually settle on its answer, and is it still computing
at the step where it appears to be reasoning? We introduce \emph{Commitment
Depth}, a causal measure that transplants one step's internal state from a
different problem and checks whether the final answer follows the transplant.
If the model has already written its answer into the state at that step, the
transplant redirects the output; if the answer has not yet been determined,
the transplant washes out. To judge whether commitment is premature, we
derive a \emph{prescribed trajectory} from the information a reference program
accumulates step by step, and prove that any faithful chain of thought must
commit no faster than this trajectory allows. We estimate Commitment Depth
with \emph{Causal Lookahead Probing} (CLP), a parameter-free intervention
that, by learning no alignment, escapes the non-linear representation dilemma
that renders learned-probe causal analyses vacuous. Across five
state-tracking tasks and four model families (Pythia, Llama-3.2, Mamba, and
the RL-tuned R1-Distill), we find that reinforcement-tuned models commit up
to $2.43\times$ earlier than the prescribed trajectory at matched accuracy,
while base and instruction-tuned models respect the bound. On
Tracr-compiled models with known computation, CLP recovers the prescribed
trajectory within its confidence interval. These results provide the first
causal evidence that reinforcement-tuned chains of thought often replay a
settled answer rather than derive it.\footnote{Code, compiled programs, and data: \texttt{\small[anonymous URL]}.}
\end{abstract}

% =====================================================================
%  1  INTRODUCTION  (related work folded in)
% =====================================================================
\section{Introduction}

A chain of thought (CoT) is not a stylistic flourish; it changes what a transformer
can compute. A constant-depth model is confined to $\mathsf{TC}^0$; the same
architecture, generating intermediate tokens, climbs to $\mathsf{P}$, each
token contributing a serial state unit the fixed-depth network lacks
\citep{DBLP:conf/iclr/MerrillS24,DBLP:conf/nips/FengZGY0W23,DBLP:journals/corr/abs-2112-00114}.
This is now practice: reinforcement-learned reasoning models rely on long chains
of thought \citep{DBLP:journals/nature/GuoYZSWZXZMBZY025}, and their traces are
read as decision evidence by humans and downstream models alike
\citep{DBLP:journals/corr/abs-2507-11473}. If the serial computation lives in
the chain, that is where we should be able to watch it happen.

Whether we can is unclear, since the expressivity theorems guarantee only that a
faithful execution \emph{exists}, not that a trained model's trace \emph{is}
one. Transformers solve algorithmic tasks just as well when reasoning tokens are
replaced by filler supplying only serial positions
\citep{DBLP:journals/corr/abs-2404-15758,DBLP:conf/iclr/GoyalJRMKN24}; and even
when tokens do carry content, they are not always operative: models adopt
answers from unverbalized features \citep{DBLP:journals/corr/abs-2503-08679} and rationalize
settled conclusions \citep{DBLP:journals/corr/abs-2505-05410}. Misleading
prompts steer chains to wrong answers even when the model would otherwise
succeed \citep{DBLP:conf/nips/TurpinMPB23}.

These observations have motivated a sustained effort to measure faithfulness
from inputs and outputs alone.
\citet{DBLP:conf/acl/JacoviG20} gave it a formal counterfactual criterion.
Subsequent work operationalized it through chain truncation
\citep{DBLP:journals/corr/abs-2307-13702}, step corruption
\citep{DBLP:journals/corr/abs-2602-11201}, and parameter erasure
\citep{DBLP:conf/emnlp/TutekCMB25}. The verdict is sobering, but the evidence
rests on behavioral signals that are themselves fragile:
\citet{DBLP:journals/tmlr/BenthamSM24} showed that a leading metric largely
re-encodes accuracy and answer-letter bias. Behavioral tests can detect that
something is wrong; they cannot locate what, inside the model, is wrong.

A second body of work looks inside and finds the answer legible early. The logit
lens \citep{nostalgebraist2020logitlens}, tuned lens
\citep{DBLP:journals/corr/abs-2303-08112}, and forward lens
\citep{DBLP:conf/conll/PalSYWB23} reveal the eventual answer from the first few
tokens; probing reasoning traces directly confirms this
\citep{DBLP:journals/corr/abs-2603-05488,
DBLP:journals/corr/abs-2603-01437,DBLP:journals/corr/abs-2511-14773,
DBLP:journals/corr/abs-2604-06613}. The temptation is to conclude that the model
has committed, but the inference fails: a linearly recoverable direction need
not play any causal role in the output
\citep{DBLP:conf/nips/VigGBQNSS20,DBLP:conf/nips/GeigerLIP21}. Commitment
demands a counterfactual: \emph{altering} the state must change the answer, not
merely \emph{reading} it.

The causal toolkit supplies such counterfactuals. Mediation analysis, activation
patching, and circuit discovery locate causally relevant components
\citep{DBLP:conf/emnlp/StolfoBS23,DBLP:conf/nips/MengBAB22,DBLP:journals/corr/abs-2404-15255,DBLP:conf/iclr/WangVCSS23},
and causal abstraction formalizes whether a network implements a hypothesized
algorithm \citep{DBLP:journals/jmlr/GeigerIZCC0AWGP25}. The sharpest tool in
this family, distributed alignment search (DAS), learns a rotation $R\!\in\!O(d)$ so that
subspace interventions match interventions on a hypothesized variable
\citep{DBLP:conf/clear2/GeigerWPIG24}. The flexibility of this rotation is also
its weakness. \citet{sutter2025nonlinear} proved that a sufficiently expressive
alignment maps any algorithm onto any network, reaching perfect scores even on
randomly initialized models, so a high alignment score certifies nothing.
Restricting to linear maps is one escape, but linearity is contested
\citep{DBLP:conf/icml/ParkCV24,DBLP:conf/iclr/BurnsYKS23} and the choice
silently determines the verdict. This is the first gap: a causal test of
commitment that rests on no learned alignment at all is missing.

A second gap is logically prior. Even with a trustworthy measure of
\emph{when} a model commits, we lack a principled account of when it
\emph{should}. Earliness is suspect only relative to the task's information
demands: if the answer is genuinely determined after a few symbols, committing
then is correct, and calling it premature would be a false alarm. The
\textsc{Rasp} family compiles readable algorithms into transformer
weights \citep{DBLP:conf/icml/WeissGY21}, and its
counting variant \textsc{C-Rasp} is the tightest known expressivity bound for
future-masked transformers, with a depth hierarchy over serial cost
\citep{DBLP:conf/nips/YangCCN25}. A $\CRASPCoT$ program states, symbol by
symbol, when the answer becomes determined and thus when commitment is licensed.
No prior work has used such a program to set a normative bar on the dynamics of
an emitted CoT.

We close both gaps at once. We define \textbf{Commitment Depth}
$\CD(\stepk)$: transplant the step-$\stepk$ residual stream from a different
problem and observe how often the final answer follows. If the model has already
written its answer into the state at step $\stepk$, the transplant redirects the
output; if the answer has not yet been determined there, the transplant washes
out. We pair this measure with a \textbf{prescribed trajectory}
$\CDP(\stepk)$ that a $\CRASPCoT$ reference program licenses, prove a
faithfulness bound $\CDM(\stepk)\le\CDP(\stepk)$, and estimate $\CD$ with
\textbf{Causal Lookahead Probing} (\CLP), a parameter-free intervention that
learns no alignment and therefore lies outside the construction that makes
alignment search vacuous. The cost is that \CLP\ measures \emph{whether} a
model has committed, not \emph{which variable} carries the commitment. For the
question of this paper, the former is what matters. Our main contributions are
as follows.
\begin{itemize}[leftmargin=*,topsep=2pt,itemsep=2pt]

\item We prove that any correct executor of a state-tracking task must satisfy
$\CDM(\stepk)\le\CDP(\stepk)$, a bound that is task-specific but
algorithm-independent, turning premature commitment into a falsifiable claim
(\cref{sec:cd}; proof in \cref{app:proof}).

\item \CLP\ is, to our knowledge, the first causal measure of CoT commitment
timing that requires no learned alignment, placing it outside the non-linear
representation dilemma of \citet{sutter2025nonlinear} (\cref{sec:clp};
\cref{app:dilemma}).

\item We compile $\CRASPCoT$ programs for five tasks spanning bounded-depth
bracket matching, solvable and non-solvable group computation, finite-state
execution, and natural-language entity tracking, and validate \CLP\ by
recovering $\CDP$ on Tracr-compiled models within its confidence interval
(\cref{sec:setup,sec:results}).

\item At matched final-answer accuracy, the RL-tuned \Rdistill\ violates the
bound on every task, committing up to $2.43\times$ earlier than $\CDP$, while
base, instruction-tuned, and state-space models all respect it ($U<0.035$).
This is the first causal evidence that reinforcement-tuned CoT often replays a
settled answer rather than derives it (\cref{sec:results}).

\end{itemize}

% =====================================================================
%  2  PRELIMINARIES AND PROBLEM SETUP
% =====================================================================
\section{Preliminaries and Problem Setup}
\label{sec:prelim}

We study CoT traces as computational objects. This section introduces the
notation for models and their intermediate states, the programming language we
use to specify a reference algorithm for each task, and the causal notion of
faithfulness that the rest of the paper builds on.

\subsection{Chain-of-Thought Computation}
\label{sec:prelim-cot}

Let $\mathcal{M}$ be an autoregressive language model over a finite vocabulary
$\Sigma$, with $L$ transformer layers and model dimension $d$. Each layer reads
and updates a hidden state vector at every position; we call this per-position
vector the \emph{residual stream} (the running sum along the residual
connections) and write $\mathbf{h}^{(\ell)}_i \in \mathbb{R}^{d}$ for it at
layer $\ell$ and position $i$.

A \emph{chain of thought} on a problem instance is the sequence of tokens
$x_{1}, x_{2}, \ldots, x_{T}$ that $\mathcal{M}$ generates after the prompt and
before the answer is read, where $T$ is the length of the trace and the
\emph{step index} $\stepk \in \{1, \ldots, T\}$ counts generated tokens. We
write $x_{1{:}\stepk} = x_{1}, \ldots, x_{\stepk}$ for the prefix through step
$\stepk$. The model's \emph{answer} is a discrete label $a \in \mathcal{A}$
obtained by applying a fixed readout $\kappa$ to the logits at the final
position, $a = \kappa(\mathcal{M}(x_{1{:}T}))$; for the state-tracking tasks we
study, $\mathcal{A}$ is the finite set of valid answers (group elements, binary
labels, automaton states, or entity locations) and $\kappa$ returns the
highest-scoring one.

To probe commitment we adopt the standard interpretability convention of
\emph{interchange interventions}: we pair a \emph{base} instance $\mathbf{b}$
(the problem under study) with a \emph{source} instance $\mathbf{s}$ (a
different problem whose gold answer differs), run $\mathcal{M}$ on both, and at
a chosen step $\stepk$ replace the base run's residual stream with the source
run's. If the final answer changes to follow the source, the model's state at
step $\stepk$ was still causally relevant; if the answer stays with the base,
the model had already committed. We use $\mathbf{b}$ and $\mathbf{s}$ throughout
in this sense.

\subsection{Prescribing an Algorithm: \texorpdfstring{$\CRASPCoT$}{C-RASP-CoT}}
\label{sec:prelim-crasp}

To judge whether commitment is premature, we need a reference algorithm whose
information demands are explicit: one that states, step by step, how much it has
learned about the answer. We take this reference from \textsc{C-Rasp}, a
programming language designed to characterize what transformers can compute
\citep{DBLP:conf/icml/WeissGY21,
DBLP:journals/corr/abs-2404-04393,DBLP:conf/nips/YangCCN25}. A \textsc{C-Rasp}
program operates on an input string $w = w_{1}\cdots w_{n} \in \Sigma^{+}$; its
key primitive, \emph{prefix counting}, lets it track quantities such as bracket
nesting depth or a running group product, and the nesting depth of these counts
measures the serial computation the program requires
\citep{DBLP:conf/nips/YangCCN25}. A \textsc{C-Rasp} program can be compiled
into transformer weights via Tracr \citep{DBLP:conf/nips/LindnerKFRMM23},
producing a model that provably executes it; we use this for validation. The
full primitive set is in \cref{app:craspcot-standard}.

A standard \textsc{C-Rasp} program reads its entire input at once, but a CoT
emits one symbol at a time and conditions on what it has already written. We
therefore use $\CRASPCoT$, the autoregressive reading of a \textsc{C-Rasp}
program in which the program consumes the prefix $x_{1{:}\stepk}$ available at
step $\stepk$ and may read its own previous outputs as input to later steps.
The semantics and equivalence to standard \textsc{C-Rasp} under this reading are
given in \cref{app:craspcot}; for the present section the program is a map
$\pi$ from prefixes to a distribution over the answer set $\mathcal{A}$.

The reference a program $\pi$ provides is informational. Let $\mathcal{D}$ be
the task distribution over instances, let $a$ be the answer as a random
variable, and let $H(\cdot)$ denote Shannon entropy in bits. The
\emph{information accumulated} by $\pi$ through step $\stepk$ is the reduction
in the answer's uncertainty given the prefix the program has read,
\begin{equation}
\label{eq:infoacc}
I_{\pi}(\stepk) \;=\; H(a) \;-\; H\!\left(a \mid x_{1{:}\stepk}\right),
\end{equation}
which increases from $0$ at $\stepk = 0$ to $H(a)$ once the prefix determines
the answer. A program that reveals the answer early has a fast-rising
$I_{\pi}$; one that requires reading the entire input has a slowly rising one.
This quantity is the basis from which we derive, in the next section, the
trajectory that a faithful execution may not outrun.

\subsection{Faithfulness as a Causal Property}
\label{sec:prelim-faithfulness}

A CoT is \emph{faithful} when the steps it emits are the steps the model uses
to reach its answer \citep{DBLP:conf/acl/JacoviG20}. This is a statement about
the model's computation, not about the plausibility or correctness of the text.
As discussed in the introduction, behavioral tests (truncating the chain,
corrupting a step) and representational tests (probing hidden states) each fall
short: the former are confounded by accuracy and format bias, the latter
establish decodability but not causation.

We therefore adopt a causal reading. A CoT has \emph{committed} to an answer at
step $\stepk$ if the model's internal state at that step determines the answer:
replacing it with the state from a different problem redirects the output.
Conversely, if such a replacement leaves the answer unchanged, the answer has
\emph{not} been committed at step $\stepk$, because it is derived from
computation that comes later. The next section formalizes this as Commitment
Depth, a continuous quantity that can be measured at every step and compared
against the prescribed trajectory derived from $I_{\pi}$.

% =====================================================================
%  3  COMMITMENT DEPTH AND THE PRESCRIBED TRAJECTORY
% =====================================================================
\section{Commitment Depth and the Prescribed Trajectory}
\label{sec:cd}

The causal reading of faithfulness in \cref{sec:prelim-faithfulness} is a
yes-or-no test at a single step. We now turn it into two continuous
trajectories: one recording when a model commits, one stating when committing is
warranted.

\subsection{Commitment Depth}
\label{sec:cd-def}

Fix a model $\mathcal{M}$ and a base--source pair $(\mathbf{b},\mathbf{s})$
with $a(\mathbf{b})\neq a(\mathbf{s})$ as defined in \cref{sec:prelim-cot}.
The requirement that the two answers differ makes the intervention
\emph{impactful}: a flip is observable only when there is a different answer to
flip to \citep{DBLP:conf/nips/GeigerLIP21}.

\begin{definition}[Commitment Depth]
\label{def:cd}
Run $\mathcal{M}$ on $\mathbf{b}$. At CoT step $\stepk$, replace the
last-layer residual stream $\mathbf{h}^{(L)}$ at the step-$\stepk$ position with
the value it takes when $\mathcal{M}$ is run on $\mathbf{s}$, then let
$\mathcal{M}$ complete the remaining chain autoregressively and read the answer
through $\kappa$. The \emph{Commitment Depth} of $\mathcal{M}$ at step $\stepk$
is the probability, over base--source pairs drawn from the task distribution
$\mathcal{D}$ with $a(\mathbf{b})\neq a(\mathbf{s})$, that this intervention
drives the final answer to the source answer:
\begin{equation}
\label{eq:cd}
\CDM(\stepk) \;=\; \Pr_{(\mathbf{b},\mathbf{s})}\!\left[\,
\kappa\!\left(\mathcal{M}^{\,\mathbf{h}^{(L)}_{\stepk}\leftarrow \mathbf{s}}(\mathbf{b})\right)
= a(\mathbf{s})\,\right].
\end{equation}
\end{definition}

The quantity reads directly as a measure of how far the model's answer has
been written into the residual stream by step $\stepk$. When $\CDM(\stepk)$ is
high, the step-$\stepk$ state carries the answer: transplanting a different
state redirects the output, confirming that the model has committed by that
point. When $\CDM(\stepk)$ is low, the answer has not yet been encoded at step
$\stepk$: the transplant has no effect because the model derives its answer from
subsequent computation. The contrast is sharp in practice: on $A_5/S_5$ at
normalized step $\stepk=0.35$, the RL-tuned model's $\CDM$ already exceeds
$0.60$ while the prescribed trajectory $\CDP$ is below $0.04$, a gap that
constitutes a bound violation. The intervention is the entire last-layer
residual stream at a single step, and it is parameter-free; the reasons for this
choice, and its relation to other interventions, are given in \cref{sec:clp}.

\subsection{The Prescribed Trajectory}
\label{sec:cd-prescribed}

Commitment Depth records when a model commits, but earliness is only suspect
relative to the information the task makes available. We therefore set against
it a normative trajectory derived from the $\CRASPCoT$ program $\pi$ for the
task. Recall from \cref{eq:infoacc} that $I_{\pi}(\stepk)=H(a)-H(a\mid
x_{1{:}\stepk})$ is the information the program has accumulated about the answer
by step $\stepk$.

\begin{definition}[Prescribed trajectory]
\label{def:cdp}
The \emph{prescribed trajectory} is $\CDP(\stepk)=g\!\left(I_{\pi}(\stepk)\right)$,
where the link function $g:[0,H(a)]\to[0,1]$ is non-decreasing with $g(0)=0$ and
$g(H(a))=1$.
\end{definition}

The link function translates accumulated information into the largest
commitment a faithful executor is entitled to: it answers the question ``given
that the program has resolved this much uncertainty about the answer, how high
can $\CDM$ be without violating faithfulness?'' Its explicit form, derived from
Fano's inequality in \cref{app:proof}, is $g(\iota)=\iota/H(a)$, so
$\CDP(\stepk)=I_{\pi}(\stepk)/H(a)$: the prescribed trajectory is simply the
fraction of the answer's uncertainty that the program has removed by step
$\stepk$.

\subsection{The Faithfulness Bound}
\label{sec:cd-bound}

The measured and prescribed trajectories are linked by a single inequality.

\begin{theorem}[Faithfulness bound]
\label{thm:bound}
If $\mathcal{M}$ faithfully executes $\pi$ in the causal sense of
\cref{sec:prelim-faithfulness}, then $\CDM(\stepk)\le\CDP(\stepk)$ for every
step $\stepk$.
\end{theorem}

\noindent The proof, which formalizes faithful execution as a constructive
abstraction and derives $g$ via Fano's inequality, is given in
\cref{app:proof}.

The bound is one-directional and, importantly, it lower-bounds \emph{any}
correct executor of the task rather than fingerprinting one particular
algorithm. Trained models are known to realize idiosyncratic and non-canonical
solutions: networks trained on the same task settle into qualitatively
different algorithms depending on initialization
\citep{DBLP:conf/nips/ZhongLTA23}, and language models trained to track state
adopt one of several distinct mechanisms \citep{DBLP:conf/icml/LiGA25}. Because
$\CDP$ constrains all of them alike, a measured violation
$\CDM(\stepk)>\CDP(\stepk)$ is diagnostic of premature commitment whatever
algorithm the model happens to implement.

\subsection{Quantifying Premature Commitment}
\label{sec:cd-functional}

A violation is summarized by how far, and how early, the model outruns its
evidence.

\begin{definition}[Unfaithfulness functional and commitment step]
\label{def:u}
The \emph{unfaithfulness functional} is the largest signed gap by which a model
exceeds the bound,
\begin{equation}
\label{eq:u}
U(\mathcal{M},\pi) \;=\; \max_{\stepk}\,\bigl[\CDM(\stepk)-\CDP(\stepk)\bigr]_{+},
\end{equation}
and the \emph{commitment step} is the first step at which the model is more
likely than not to have committed, $\stepk^{*}=\min\{\stepk:\CDM(\stepk)\ge
\tfrac{1}{2}\}$.
\end{definition}

We report $U(\mathcal{M},\pi)$ as the headline measure of how far a model
commits ahead of its evidence, and $\stepk^{*}$, which locates commitment on
the step axis, for matched-accuracy comparisons across model families.
\Cref{fig:cd-schematic} illustrates the two trajectories and the gap that $U$
measures.

\begin{figure}[t]
\centering
\includegraphics[width=0.8\columnwidth,trim=0 0 0 10,clip]{fig_cd_schematic.pdf}
\vspace{-4pt}
\caption{Commitment Depth $\CDM(\stepk)$ against the prescribed trajectory
$\CDP(\stepk)$. A faithful executor stays at or below the bound (base); a model
that commits prematurely (RL-tuned) crosses it, and $U(\mathcal{M},\pi)$ is the
largest gap. $\stepk^{*}$ marks where commitment first exceeds one half.}
\label{fig:cd-schematic}
\end{figure}

% =====================================================================
%  4  CAUSAL LOOKAHEAD PROBING
% =====================================================================
\section{Causal Lookahead Probing}
\label{sec:clp}

Causal Lookahead Probing (\CLP) is the estimator of $\CDM(\stepk)$: it
transplants a single last-layer state from a counterfactual run, lets the model
finish its own chain, and records whether the answer follows. The design is
deliberately rigid, and that rigidity is what makes its readings causally
interpretable.

\subsection{The Intervention}
\label{sec:clp-intervention}

\Cref{alg:clp} states the procedure. For each base--source pair, we cache the
source's last-layer state at step $\stepk$, inject it into the base run at the
same step, let the model finish autoregressively, and record whether the answer
follows the source. Averaging $N$ such indicators yields $\widehat{\CDM}
(\stepk)$; the sampling design and variance analysis are in
\cref{app:metrics}.

\begin{algorithm}[t]
\caption{Causal Lookahead Probing at step $\stepk$}
\label{alg:clp}
\begin{algorithmic}[1]
\Require model $\mathcal{M}$, step $\stepk$, pairs $\{(\mathbf{b}_n,\mathbf{s}_n)\}_{n=1}^{N}$ with $a(\mathbf{b}_n)\neq a(\mathbf{s}_n)$
\State $c \gets 0$
\For{$n = 1$ to $N$}
  \State run $\mathcal{M}$ on $\mathbf{s}_n$; cache $\mathbf{h}^{(L)}_{\stepk}$
  \State run $\mathcal{M}$ on $\mathbf{b}_n$; at step $\stepk$ set $\mathbf{h}^{(L)}_{\stepk}\!\gets$ cached value
  \State complete the chain autoregressively; read $\hat{a}=\kappa(\cdot)$
  \If{$\hat{a} = a(\mathbf{s}_n)$} $c \gets c + 1$ \EndIf
\EndFor
\State \Return $\widehat{\CDM}(\stepk) = c / N$
\end{algorithmic}
\end{algorithm}

The intervention transplants the entire last-layer residual stream at one step
and leaves every parameter untouched; it is, in the language of
\cref{sec:prelim-faithfulness}, a counterfactual on the model's own state rather
than a readout of it.

\subsection{Immunity to the Non-Linear Representation Dilemma}
\label{sec:clp-immunity}

Learned causal analyses fit an alignment between a hypothesized variable and the
network's activations. As discussed in the introduction, the flexibility of the
learned map is the source of the vacuity problem: a sufficiently expressive
alignment achieves perfect interchange accuracy even on randomly initialized
models \citep{sutter2025nonlinear}. \CLP\ sidesteps this entirely. Its alignment
is the identity on the full residual stream, $\tau = \mathrm{id}$: no rotation
$R$, no subspace rank $r$, and no optimization of any kind.

\begin{proposition}[Immunity to the dilemma]
\label{prop:immunity}
Because \CLP\ fixes the alignment to the identity on the full last-layer
residual stream and optimizes no alignment parameters, the family of expressive
maps that drives the non-linear representation dilemma cannot be applied to it.
A nonzero or bound-violating reading of $\CDM(\stepk)$ therefore cannot be
manufactured by alignment flexibility.
\end{proposition}

\noindent The formal statement, including the class of maps the dilemma quantifies
over and why $\tau=\mathrm{id}$ lies outside it, is given in \cref{app:dilemma}.

The relation of \CLP\ to activation patching, distributed alignment search,
layerwise causal mediation, and forward-decoding probes is detailed in
\cref{sec:clp-relation}; in each case the distinction is the same: the existing
method asks where information is or whether it can be read out, while \CLP\ asks
whether altering it moves the answer.

\subsection{Scope and Trade-offs}
\label{sec:clp-scope}

We are explicit about what \CLP\ does not deliver. By transplanting the entire
last-layer state, it forgoes variable-level attribution: it measures whether the
model has committed to an answer, not which internal variable carries that
commitment. We accept this knowingly, because the question of this paper is
about the answer, when a model fixes it relative to when the information
licenses fixing it, and answering it does not require localizing the variable
that holds the commitment; that localization is left to future work.

The choice of the last layer follows from the same logic. The last-layer
residual stream feeds the unembedding directly, so it is the state closest to
the emitted answer and the one in which answer-relevant information is most fully
integrated; mechanistic studies of arithmetic find that the result of a
computation is written into the residual stream by the late feed-forward layers
\citep{DBLP:conf/emnlp/StolfoBS23}, consistent with reading commitment there. The
sensitivity of \CLP\ to this choice, including measurements at earlier layers, is
examined in the ablation suite (\cref{app:ablations}).

% =====================================================================
%  5  EXPERIMENTAL SETUP
% =====================================================================
\section{Experimental Setup}
\label{sec:setup}

We choose tasks for which a $\CRASPCoT$ program, and hence a prescribed
trajectory, can be written down, and models that let us hold final-answer
accuracy fixed while varying whether a model was trained to reason. Both
choices are what give a bound violation its force.

\subsection{Tasks}
\label{sec:setup-tasks}

The five tasks are state-tracking problems ordered by the algebraic difficulty
that determines when their answer becomes determined, and therefore by how
diagnostic premature commitment is. \textsc{Dyck-2}, balanced bracketing with
bounded nesting depth, fixes its answer early, so its prescribed trajectory
rises soon and it serves mainly to confirm that the estimator does not report
spurious violations when the task leaves little room for one
\citep{DBLP:conf/emnlp/EbrahimiGZ20}.
The group word problems compute a product in a finite group from a sequence of
generators: the solvable groups $A_4$ and $S_4$ admit shallow $\CRASPCoT$
solutions, whereas the non-solvable groups $A_5$ and $S_5$ are
$\mathsf{NC}^1$-complete (requiring logarithmic-depth circuits that
constant-depth transformers cannot simulate without serial steps), so the answer
stays genuinely uncertain until late. The non-solvable case is therefore the
sharpest test: its prescribed trajectory rises slowly, and any model that
commits early must outrun it. Two further tasks supply ecological contrast:
finite-state automaton execution, where the model is given a DFA and an input
string and must report the final state
\citep{DBLP:conf/acl/ZhangDJ0J25}, and natural-language entity tracking, where
the model reads a short discourse describing objects changing location and
reports the final position of a queried entity. Both reduce to the same
state-tracking template \citep{DBLP:conf/icml/LiGA25}; construction details
are in \cref{app:exp}. \Cref{tab:tasks} summarizes all five.

\begin{table}[t]
\centering
\small
\setlength{\tabcolsep}{4pt}
\begin{tabular}{@{}llccc@{}}
\toprule
Task & Class & $|\mathcal{A}|$ & $\CDP$ rise & GT \\
\midrule
\textsc{Dyck-2}      & bounded-depth      & 2  & early  & \checkmark \\
$A_4,\,S_4$          & solvable group     & 12, 24 & mid   & \checkmark \\
$A_5,\,S_5$          & non-solv.\ ($\mathsf{NC}^1$) & 60, 120 & late & \checkmark \\
FSA exec.            & regular            & varies & mid   & \checkmark \\
entity track.        & natural language   & varies & late  & --- \\
\bottomrule
\end{tabular}
\caption{The five state-tracking tasks, ordered by algebraic difficulty.
$|\mathcal{A}|$ is the answer-set size; ``$\CDP$ rise'' is when the prescribed
trajectory accumulates information; GT marks tasks with a Tracr-compiled
ground-truth model (entity tracking uses an approximate program, so no exact
GT is available).}
\label{tab:tasks}
\end{table}

\subsection{Compiled Programs and Ground Truth}
\label{sec:setup-programs}

For each task we write a $\CRASPCoT$ program $\pi$ using the Boolean and
counting primitives of \cref{sec:prelim-crasp}, which fixes the information
trajectory $I_{\pi}(\stepk)$ and hence $\CDP(\stepk)$ for that task. We then
compile $\pi$ with Tracr into a transformer of known structure
\citep{DBLP:conf/nips/LindnerKFRMM23}. Because the compiled model provably
executes $\pi$, applying \CLP\ to it should recover $\CDP$ up to estimation
error; we use this as the validity check for the method. Compilation
details, including the counting implementation and precision, appear in
\cref{app:compile}.

\subsection{Models}
\label{sec:setup-models}

We evaluate four families chosen so that reasoning training is the only axis
that varies once accuracy is matched. \Pythia\ (410M) and the base \Llama\
(3B) provide reasoning-free references; the instruction-tuned \Llama\ adds
supervised tuning without reinforcement; \Mamba\ (2.8B) is an architectural
control, since state-space models are expressively limited in much the same way
as transformers for state tracking and so should not behave differently for
architectural reasons alone \citep{DBLP:conf/icml/MerrillPS24}; and \Rdistill\
(8B) is the reinforcement-tuned reasoning model under test.

\subsection{Metrics and Protocol}
\label{sec:setup-protocol}

We compare model families within matched final-answer accuracy bins, so that
$\stepk^{*}$ is read off at equal accuracy, and report bootstrap confidence
intervals on $\widehat{\CDM}(\stepk)$. Alongside \CLP\ we place behavioral
faithfulness baselines: early-answering
\citep{DBLP:journals/corr/abs-2307-13702}, its accuracy-controlled normalization
\citep{DBLP:journals/tmlr/BenthamSM24}, and the detection--extraction gap
\citep{DBLP:journals/corr/abs-2604-06613}, together with answer and correctness
probes. Sampling design, competitor reproductions, and compute are detailed in
\cref{app:exp,app:competitors}.

% =====================================================================
%  6  RESULTS
% =====================================================================
\section{Results}
\label{sec:results}

\subsection{The bound is violated, most strongly by reinforcement-tuned models}
\label{sec:results-main}

\Cref{fig:results-main} plots $\CDM(\stepk)$ against the prescribed trajectory
$\CDP(\stepk)$ for every model family and task. The reasoning-free families,
base \Pythia\ and \Llama, the instruction-tuned \Llama, and the architectural
control \Mamba, stay close to the bound, with small unfaithfulness functionals.
The reinforcement-tuned \Rdistill\ crosses the bound, and the
violation grows with algebraic difficulty: reading $U$ across the columns of
\cref{tab:headline}, the ordering is Dyck-2 ($U\!=\!0.289$) $<$ FSA ($0.616$)
$\approx$ Entity ($0.618$) $<$ $A_4/S_4$ ($0.672$) $<$ $A_5/S_5$ ($0.748$),
consistent with the theoretical prediction that tasks requiring deeper serial
computation leave more room for premature commitment.
\Cref{tab:headline} reports $U(\mathcal{M},\pi)$ and the commitment step
$\stepk^{*}$ for all 25 (model, task) cells; the pattern is consistent across
every task: \Rdistill\ shows the largest $U$ and the earliest $\stepk^{*}$,
while all four base families stay within $U < 0.035$, none of which differs
significantly from zero. On $A_5/S_5$ the reasoning model's $U$ exceeds any
baseline by more than an order of magnitude ($p < 0.001$ for all pairwise
comparisons after Holm--Bonferroni correction; on $A_5/S_5$ specifically,
Cliff's $\delta > 0.91$; \cref{app:stattests}). Final-answer accuracy ranges from 0.62 (\Pythia\ on
$A_5/S_5$) to 0.94 (\Rdistill\ on Dyck-2), with substantial overlap between
families on every task, ensuring that the matched-accuracy comparison of the
next subsection has adequate support.

\begin{table*}[t]
\centering
\small
\setlength{\tabcolsep}{3.5pt}
\begin{tabular}{@{}l cc cc cc cc cc@{}}
\toprule
& \multicolumn{2}{c}{\textsc{Dyck-2}}
& \multicolumn{2}{c}{$A_4/S_4$}
& \multicolumn{2}{c}{$A_5/S_5$}
& \multicolumn{2}{c}{FSA}
& \multicolumn{2}{c}{Entity} \\
\cmidrule(lr){2-3}\cmidrule(lr){4-5}\cmidrule(lr){6-7}\cmidrule(lr){8-9}\cmidrule(lr){10-11}
Model & $U$ & $\stepk^{*}$ & $U$ & $\stepk^{*}$ & $U$ & $\stepk^{*}$ & $U$ & $\stepk^{*}$ & $U$ & $\stepk^{*}$ \\
\midrule
\Pythia        & 0.016 & 0.272 & 0.004 & 0.485 & 0.034 & 0.793 & 0.025 & 0.564 & 0.011 & 0.704 \\
\Llama\ (base) & 0.005 & 0.268 & 0.006 & 0.486 & 0.008 & 0.751 & 0.010 & 0.536 & 0.009 & 0.681 \\
\Llama\ (inst) & 0.024 & 0.263 & 0.012 & 0.476 & 0.011 & 0.742 & 0.016 & 0.534 & 0.012 & 0.706 \\
\Mamba         & 0.014 & 0.267 & 0.008 & 0.476 & 0.011 & 0.741 & 0.005 & 0.542 & 0.026 & 0.690 \\
\Rdistill      & \textbf{0.289}$^{\dagger}$ & \textbf{0.178} & \textbf{0.672}$^{\dagger}$ & \textbf{0.187} & \textbf{0.748}$^{\dagger}$ & \textbf{0.304} & \textbf{0.616}$^{\dagger}$ & \textbf{0.274} & \textbf{0.618}$^{\dagger}$ & \textbf{0.361} \\
\bottomrule
\end{tabular}
\caption{Unfaithfulness functional $U$ and commitment step $\stepk^{*}$ across
all five tasks and five models. $^{\dagger}$Significantly greater than all base
models after Holm--Bonferroni correction ($p<0.001$, permutation test with
$10{,}000$ resamples; full results in \cref{app:stattests}). Bootstrap 95\% CI
widths are $\le 0.04$ for $U$ and $\le 0.03$ for $\stepk^{*}$; no base model's
$U$ differs significantly from zero.}
\label{tab:headline}
\end{table*}

\subsection{At matched accuracy, reinforcement tuned models commit earlier}
\label{sec:results-matched}

A natural objection is that earlier commitment merely reflects greater
capability. \Cref{fig:matched-acc} controls for this by comparing the
commitment step $\stepk^{*}$ within matched final-answer accuracy bins (five
equal-width bins from 0.6 to 1.0, each containing at least 50 instances). At
equal accuracy on $A_5/S_5$, \Rdistill\ commits at a substantially smaller
$\stepk^{*}$ than the base families and than the prescribed trajectory, a gap of
$2.43\times$ on the step axis in the highest accuracy bin (bootstrap 95\% CI:
$[2.07, 2.91]$). Because accuracy is held fixed, the earliness cannot be
attributed to the reasoning model simply being better at the task. The same
pattern holds on all five tasks: within each accuracy bin, \Rdistill\ commits
earlier than every base family (\cref{app:perdataset}).

\subsection{Validity: \CLP\ recovers the prescribed trajectory}
\label{sec:results-validity}

\Cref{fig:validity} checks that \CLP\ measures what it claims. On
Tracr-compiled models that provably execute the task program, the estimated
$\CDM(\stepk)$ tracks $\CDP(\stepk)$ within the bootstrap confidence interval
at every step, with a maximum absolute deviation of 0.023 on $A_5/S_5$ and
below 0.02 on all four exact-program tasks.
A control that injects a random or orthogonal residual state in place of the
source state moves the answer at chance rate ($0.498 \pm 0.021$), confirming that
the flips counted by \CLP\ reflect the transplanted computation rather than a
generic perturbation artifact.

\subsection{Robustness}
\label{sec:results-robustness}

Four ablation checks support the reading above (full panels in
\cref{app:ablations,app:perdataset,app:competitors}). \emph{Layer ablation}:
$\CDM$ at layer $L/2$ drops by $49\% \pm 7\%$ relative to layer $L$ on
$A_5/S_5$, and early layers read at chance ($0.51 \pm 0.03$). \emph{Control
injection}: random and orthogonalized state replacements yield flip rates of
$0.503 \pm 0.018$, confirming that measured commitment reflects transplanted
computation; same-answer sources produce rates below $0.017 \pm 0.008$.
\emph{Decodability precedes commitment}: linear probes reach 80\% accuracy
$0.147 \pm 0.021$ normalized steps before $\CDM$ reaches 0.5 on \Rdistill,
confirming the gap between decodability and causal commitment.
\emph{Source and temperature stability}: varying source selection shifts
$\stepk^{*}$ by at most $0.028 \pm 0.011$; varying temperature from 0 to 1
shifts it by at most $0.047 \pm 0.019$, with model ordering preserved
throughout.


\begin{figure}[t]
\centering
\includegraphics[width=0.9\columnwidth]{fig_results_compact.pdf}

\caption{Commitment Depth $\CDM(\stepk)$ against the prescribed trajectory
$\CDP(\stepk)$ (dashed) on the non-solvable group $A_5/S_5$. Only \Rdistill\
crosses the bound; the full model$\times$task grid is in \cref{app:perdataset}.}
\label{fig:results-main}
\end{figure}

\begin{figure}[t]
\centering
\includegraphics[width=0.85\columnwidth]{fig_matched_acc.pdf}

\caption{Commitment step $\stepk^{*}$ within matched final-answer accuracy bins
on $A_5/S_5$. \Rdistill\ commits earlier than the base families at equal
accuracy.}
\label{fig:matched-acc}
\end{figure}

\begin{figure}[t]
\centering
\includegraphics[width=0.85\columnwidth]{fig_validity.pdf}

\caption{On a Tracr-compiled ground-truth model, \CLP\ recovers $\CDP$, while
injecting a random residual state moves the answer at chance.}
\label{fig:validity}
\end{figure}

% =====================================================================
%  7  ANALYSIS, DISCUSSION, LIMITATIONS
% =====================================================================
\section{Discussion}
\label{sec:discussion}

What does the gap between $\CDM$ and $\CDP$ tell us about what RL training does
to the model's computation? Reinforcement on outcome correctness rewards
reaching the right answer by any means, including means that bypass the serial
reasoning the CoT is supposed to externalize. The model learns to commit early
because early commitment followed by plausible filler is rewarded just as much
as genuine derivation. The CoT then functions less as a workspace and more as a
transcript of a decision already made, consistent with behavioral findings that
reasoning models adopt answers from unverbalized features and rationalize
settled conclusions \citep{DBLP:journals/corr/abs-2503-08679}.
Our contribution is to make this observation causal and normatively grounded:
the bound says not only that the model committed early, but that it committed
earlier than any correct executor is entitled to.

This finding has practical consequences. If CoT traces are used as explanations
or as safety-relevant monitoring signals, premature commitment means the
explanation does not reflect the computation it purports to describe. A
practitioner reading the trace of an RL-tuned model should not assume that
the reasoning steps caused the answer; our measurements suggest that, on
state-tracking tasks, the answer was often fixed before the critical reasoning
steps were emitted. Whether this extends to open-ended reasoning, where no
$\CRASPCoT$ reference is available, is an open question, but the implication
for controlled tasks is clear: the chain is not always the computation.

Natural extensions include using $\CDM$ as a training signal to penalize
premature commitment during RL fine-tuning, tracking $\CDM$ across training
checkpoints to locate when premature commitment first emerges, and extending the
prescribed trajectory beyond $\CRASPCoT$ to bring open-ended reasoning tasks
within scope.

% ===== Limitations (required ACL section) =====
\section*{Limitations}

Our prescribed trajectory is defined for tasks that admit a $\CRASPCoT$ program,
which covers the state-tracking problems we study but not open-ended reasoning,
where no such reference algorithm is available; extending the framework beyond
this class is open. Commitment Depth is estimated at the last-layer residual
stream, so it captures when an answer becomes causally fixed but not which
intermediate variable carries it, and a finer-grained account would require the
learned-alignment machinery we deliberately avoid. The natural-language entity
tracking task is the least controlled of the five, since its prescribed
trajectory is approximate rather than exact. Our model families span base,
instruction-tuned, state-space, and reinforcement-tuned systems, but a single
reinforcement-tuned family cannot establish that the effect holds across all
training recipes, and broader coverage is needed to claim generality.

% ===== Ethics (optional ACL section) =====
\section*{Ethics Statement}

This work studies whether the CoT traces emitted by reasoning models
reflect the computation behind their answers. The question bears directly on the
practice of reading such traces as explanations or as monitoring signals: a
trace that performs a settled answer rather than computing it can mislead a
reader who treats it as faithful. By giving a causal and normative handle on
when a model commits, our method is intended to make that practice more
cautious, not to enable any new capability. The models and tasks we use are
standard and involve no human subjects or sensitive data.

% =====================================================================
%  BIBLIOGRAPHY
% =====================================================================
\bibliography{emnlp}

% =====================================================================
%  APPENDIX
% =====================================================================
\appendix

\section*{Appendix Table of Contents}
\vspace{-3pt}
{\normalsize\setlength{\leftmargini}{1.8em}
\begin{itemize}[itemsep=4pt,topsep=2pt,leftmargin=*]
\item[\textbf{A}] \hyperref[app:notation]{\textbf{Notation}}
\item[\textbf{B}] \hyperref[app:proof]{\textbf{Proof of the Faithfulness Bound}}\\[1pt]
  \hyperref[app:proof-abstraction]{B.1 Faithful Execution as a Constructive Abstraction}\\
  \hyperref[app:proof-cd]{B.2 The Commitment Depth Intervention in Terms of Information}\\
  \hyperref[app:proof-g]{B.3 Construction of the Link Function}\\
  \hyperref[app:proof-main]{B.4 Proof of Theorem 1}\\
  \hyperref[app:proof-boundary]{B.5 Boundary Cases and Scope of the Bound}
\item[\textbf{C}] \hyperref[app:dilemma]{\textbf{CLP and the Non-Linear Representation Dilemma}}\\[1pt]
  \hyperref[app:dilemma-statement]{C.1 The Dilemma, Formally}\\
  \hyperref[app:dilemma-clp]{C.2 CLP as a Degenerate Alignment}\\
  \hyperref[app:dilemma-proof]{C.3 Proof of Proposition 1}\\
  \hyperref[app:dilemma-others]{C.4 Dilemma Status of Other Causal Methods}\\
  \hyperref[sec:clp-relation]{C.5 Relation to Existing Interventions}
\item[\textbf{D}] \hyperref[app:craspcot]{\textbf{C-RASP\textsuperscript{CoT} Semantics and Equivalence}}\\[1pt]
  \hyperref[app:craspcot-standard]{D.1 Standard C-RASP Semantics}\\
  \hyperref[app:craspcot-auto]{D.2 The Autoregressive Reading}\\
  \hyperref[app:craspcot-equiv]{D.3 Equivalence}
\item[\textbf{E}] \hyperref[app:compile]{\textbf{Compilation Pipeline}}\\[1pt]
  \hyperref[app:compile-overview]{E.1 Overview}\\
  \hyperref[app:compile-pertask]{E.2 Per-Task Programs}\\
  \hyperref[app:compile-params]{E.3 Compilation Parameters}\\
  \hyperref[app:compile-precision]{E.4 Precision Analysis}
\item[\textbf{F}] \hyperref[app:exp]{\textbf{Experimental Details}}\\[1pt]
  \hyperref[app:exp-models]{F.1 Models}\\
  \hyperref[app:exp-tasks]{F.2 Tasks}\\
  \hyperref[app:exp-sampling]{F.3 Sampling Design}\\
  \hyperref[app:exp-generation]{F.4 Generation Configuration}\\
  \hyperref[app:exp-compute]{F.5 Compute}\\
  \hyperref[app:exp-software]{F.6 Software}
\item[\textbf{G}] \hyperref[app:metrics]{\textbf{Metric Definitions}}\\[1pt]
  \hyperref[app:metrics-estimator]{G.1 The Commitment Depth Estimator}\\
  \hyperref[app:metrics-summary]{G.2 Summary Metrics}\\
  \hyperref[app:metrics-baselines]{G.3 Baseline Measures}\\
  \hyperref[app:metrics-competitors]{G.4 Competitor Method Metrics}\\
  \hyperref[app:metrics-stats]{G.5 Statistical Testing Framework}
\item[\textbf{H}] \hyperref[app:competitors]{\textbf{Competitor Reproductions}}\\[1pt]
  \hyperref[app:competitors-ea]{H.1 Early-Answering}\\
  \hyperref[app:competitors-ac]{H.2 Accuracy-Controlled Normalization}\\
  \hyperref[app:competitors-de]{H.3 Detection--Extraction Gap}\\
  \hyperref[app:competitors-das]{H.4 Distributed Alignment Search (DAS)}\\
  \hyperref[app:competitors-cmi]{H.5 Layerwise Mediation (CMI)}\\
  \hyperref[app:competitors-precot]{H.6 Pre-CoT Decoding Probes}\\
  \hyperref[app:competitors-nldd]{H.7 Step-Level Faithfulness Decay (NLDD)}\\
  \hyperref[app:competitors-overlay]{H.8 Competitor Overlay}\\
  \hyperref[app:competitors-summary]{H.9 Comprehensive Comparison}
\item[\textbf{I}] \hyperref[app:ablations]{\textbf{Full Ablation Suite}}\\[1pt]
  \hyperref[app:ablation-layer]{I.1 Layer Ablation}\\
  \hyperref[app:ablation-sampling]{I.2 Sampling Ablation}\\
  \hyperref[app:ablation-control]{I.3 Control-Injection Ablation}\\
  \hyperref[app:ablation-probes]{I.4 Probe Comparison}\\
  \hyperref[app:ablation-temperature]{I.5 Temperature Ablation}\\
  \hyperref[app:ablation-chainlen]{I.6 Chain-Length Ablation}\\
  \hyperref[app:ablation-source]{I.7 Source-Selection Ablation}\\
  \hyperref[app:ablation-readout]{I.8 Readout-Position Ablation}\\
  \hyperref[app:ablation-bootstrap]{I.9 Bootstrap Stability}\\
  \hyperref[app:ablation-summary]{I.10 Ablation Summary}
\item[\textbf{J}] \hyperref[app:perdataset]{\textbf{Per-Dataset Breakdowns}}
\item[\textbf{K}] \hyperref[app:prompts]{\textbf{Prompt Templates and Transcripts}}\\[1pt]
  \hyperref[app:prompts-group]{K.1 Group Word Problems}\\
  \hyperref[app:prompts-dyck]{K.2 Dyck-2}\\
  \hyperref[app:prompts-fsa]{K.3 Finite-State Automaton Execution}\\
  \hyperref[app:prompts-entity]{K.4 Entity Tracking}\\
  \hyperref[app:prompts-errors]{K.5 Error Examples and Commitment Signatures}
\item[\textbf{L}] \hyperref[app:negative]{\textbf{Negative Results and Failure Cases}}
\item[\textbf{M}] \hyperref[app:checklist]{\textbf{Reproducibility Checklist}}\\[1pt]
  \hyperref[app:checklist-arr]{M.1 ARR Reproducibility Checklist}\\
  \hyperref[app:checklist-hyperparams]{M.2 Hyperparameter Summary}
\item[\textbf{N}] \hyperref[app:capability]{\textbf{Method Capability Comparison}}
\item[\textbf{O}] \hyperref[app:stattests]{\textbf{Statistical Tests}}\\[1pt]
  \hyperref[app:stattests-perm]{O.1 Permutation Tests for $U$}\\
  \hyperref[app:stattests-effect]{O.2 Effect Sizes for $\stepk^{*}$}\\
  \hyperref[app:stattests-summary]{O.3 Summary}
\item[\textbf{P}] \hyperref[app:related]{\textbf{Extended Related Work}}\\[1pt]
  \hyperref[app:related-behavioral]{P.1 Behavioral Faithfulness}\\
  \hyperref[app:related-probing]{P.2 Internal-State Reading and Probing}\\
  \hyperref[app:related-causal]{P.3 Causal Intervention Methods}\\
  \hyperref[app:related-positioning]{P.4 Positioning of the Present Work}
\end{itemize}
}
\vspace{6pt}

\section{Notation}
\label{app:notation}

\Cref{tab:notation} collects the symbols used throughout the paper, grouped by
the part of the framework they belong to, with the location of each symbol's
definition. Where a symbol is introduced by a displayed definition or equation,
the reference points to it.

\tablefirsthead{\toprule Symbol & Meaning & Defined \\ \midrule}
\tablehead{\multicolumn{3}{@{}l}{\small\textit{continued}}\\ \toprule Symbol & Meaning & Defined \\ \midrule}
\tabletail{\midrule \multicolumn{3}{r@{}}{\small\textit{continued on next page}}\\}
\tablelasttail{\bottomrule}
\bottomcaption{Notation used in the paper.}
\label{tab:notation}
\begin{supertabular}{@{}p{0.16\columnwidth}p{0.56\columnwidth}p{0.18\columnwidth}@{}}
\multicolumn{3}{@{}l}{\textit{Models and chains of thought}}\\
$\mathcal{M}$ & autoregressive language model under study & \cref{sec:prelim-cot}\\
$L$, $d$ & number of transformer layers; model dimension & \cref{sec:prelim-cot}\\
$\mathbf{h}^{(\ell)}_{i}$ & residual stream at layer $\ell$, position $i$ & \cref{sec:prelim-cot}\\
$\mathbf{h}^{(L)}_{\stepk}$ & last-layer residual stream at step $\stepk$ & \cref{def:cd}\\
$\stepk$, $T$ & chain-of-thought step index; trace length & \cref{sec:prelim-cot}\\
$x_{1{:}\stepk}$ & chain-of-thought prefix through step $\stepk$ & \cref{sec:prelim-cot}\\
$\kappa$ & readout from final-position logits to an answer & \cref{sec:prelim-cot}\\
$\mathcal{A}$, $a$ & answer set; gold answer as a random variable & \cref{sec:prelim-cot}\\
$\mathbf{b}$, $\mathbf{s}$ & base and source instances of an intervention & \cref{sec:cd-def}\\
$\mathcal{D}$ & task distribution over instances & \cref{app:proof}\\
\addlinespace
\multicolumn{3}{@{}l}{\textit{Reference algorithm}}\\
$\pi$ & $\CRASPCoT$ program for the task & \cref{sec:prelim-crasp}\\
$\Sigma$, $\Gamma$ & input alphabet; output-symbol alphabet & \cref{sec:prelim-crasp}\\
$Q_{\sigma}$ & input Boolean vector for symbol $\sigma$ & \cref{app:craspcot-standard}\\
$\#P$ & prefix count of a Boolean vector $P$ & \cref{app:craspcot-standard}\\
$H(\cdot)$ & Shannon entropy in bits & \cref{sec:prelim-crasp}\\
$I_{\pi}(\stepk)$ & information about $a$ accumulated by step $\stepk$ & \cref{eq:infoacc}\\
\addlinespace
\multicolumn{3}{@{}l}{\textit{Commitment Depth and the bound}}\\
$\CDM(\stepk)$ & Commitment Depth: causal flip probability & \cref{def:cd}\\
$\CDP(\stepk)$ & prescribed trajectory $g(I_{\pi}(\stepk))$ & \cref{def:cdp}\\
$g$ & link function $[0,H(a)]\to[0,1]$ & \cref{def:cdp}\\
$U(\mathcal{M},\pi)$ & unfaithfulness functional, the largest gap & \cref{def:u}\\
$\stepk^{*}$ & commitment step, first $\stepk$ with $\CDM\ge\tfrac12$ & \cref{def:u}\\
$S_{\stepk}$, $\phi_{\stepk}$ & prefix $\sigma$-algebra; abstraction map & \cref{def:abstraction}\\
\addlinespace
\multicolumn{3}{@{}l}{\textit{Causal Lookahead Probing and alignment}}\\
$\widehat{\CDM}(\stepk)$ & Monte Carlo estimate of $\CDM(\stepk)$ & \cref{sec:clp-intervention}\\
$N$ & number of base--source pairs sampled & \cref{sec:clp-intervention}\\
$\tau$, $\Pi$ & featurizing map; coordinate selection & \cref{app:dilemma-statement}\\
$R\in O(d)$, $r$ & learned rotation and subspace rank (DAS) & \cref{sec:clp-immunity}\\
$\mathcal{F}$ & class of featurizing maps searched over & \cref{def:expressive}\\
$V$, $\mathcal{V}$ & high-level variable and its value set & \cref{app:dilemma-statement}\\
$\mathrm{IIA}(\tau,\Pi)$ & interchange-intervention accuracy & \cref{eq:iia}\\
\end{supertabular}

\section{Proof of the Faithfulness Bound}
\label{app:proof}

This appendix proves \cref{thm:bound}: if a model faithfully executes a
$\CRASPCoT$ program $\pi$, then its Commitment Depth satisfies $\CDM(\stepk)
\le\CDP(\stepk)$ at every step. The argument proceeds in four stages.

\begin{enumerate}[leftmargin=*,topsep=2pt,itemsep=1pt]
\item \textbf{Formalize faithfulness} (\cref{app:proof-abstraction}). We
  define faithful execution as a constructive abstraction: the model's eventual
  answer at step $\stepk$ depends only on the prefix $x_{1{:}\stepk}$ that $\pi$
  has read, not on any internal state beyond what the prefix determines.
\item \textbf{Bound the flip probability} (\cref{app:proof-cd}). Under this
  abstraction, transplanting a source state at step $\stepk$ can redirect the
  answer only if the source prefix has resolved the answer, which happens with
  probability at most $1-H(a\mid x_{1{:}\stepk})/H(a)$ by Fano's inequality.
\item \textbf{Construct the link function} (\cref{app:proof-g}). We define
  $g(\iota)=\iota/H(a)$, verify it satisfies the requirements of
  \cref{def:cdp}, and show that the right side of the flip bound equals
  $\CDP(\stepk)=g(I_{\pi}(\stepk))$.
\item \textbf{Combine} (\cref{app:proof-main}). Chaining the three steps yields
  $\CDM(\stepk)\le\CDP(\stepk)$.
\end{enumerate}

\noindent Throughout, the task instance is a random variable drawn from the task
distribution $\mathcal{D}$, the answer $a\in\mathcal{A}$ is a measurable
function of the instance, and $x_{1{:}\stepk}$ is the CoT prefix the model has
emitted through step $\stepk$. All entropies are in bits and all probabilities
are taken over $\mathcal{D}$ together with any sampling randomness of
$\mathcal{M}$.

\subsection{Faithful Execution as a Constructive Abstraction}
\label{app:proof-abstraction}

We make precise the causal reading of faithfulness from
\cref{sec:prelim-faithfulness}. Let $\pi$ be the $\CRASPCoT$ program for the
task and let $S_{\stepk}=\sigma\!\left(x_{1{:}\stepk}\right)$ denote the
information $\pi$ has read by step $\stepk$, formally the $\sigma$-algebra
generated by the prefix. Write $\mathcal{M}_{\stepk}$ for the conditional
distribution over the model's final answer obtained by running $\mathcal{M}$
from its state at step $\stepk$ to the end of the chain and reading $\kappa$.

\begin{definition}[Constructive abstraction]
\label{def:abstraction}
$\mathcal{M}$ \emph{faithfully executes} $\pi$ if for every step $\stepk$ the
final-answer distribution $\mathcal{M}_{\stepk}$ is measurable with respect to
$S_{\stepk}$; that is, there is a measurable map $\phi_{\stepk}$ such that
\begin{equation}
\label{eq:abstraction}
\mathcal{M}_{\stepk} \;=\; \phi_{\stepk}\!\left(x_{1{:}\stepk}\right)
\qquad\text{almost surely.}
\end{equation}
\end{definition}

\noindent\textbf{Interpretation.}
\Cref{eq:abstraction} states that, conditioned on what $\pi$ has read, the
model's eventual answer carries no further dependence on the internal state at
step $\stepk$: the prefix is a sufficient statistic for the answer the model
will give. This is the distributional form of the single-step causal criterion
of \cref{sec:prelim-faithfulness}: intervening on the state at step
$\stepk$ changes the answer only insofar as it changes the prefix the model
conditions on. The definition is stronger than behavioral faithfulness (which
tests only input-output correlations) and weaker than full causal abstraction
(which requires a variable-level alignment); it captures exactly the claim that
the model's commitment at step $\stepk$ is warranted by the prefix information.

\subsection{The Commitment Depth Intervention in Terms of Information}
\label{app:proof-cd}

Fix a base--source pair $(\mathbf{b},\mathbf{s})$ with $a(\mathbf{b})\neq
a(\mathbf{s})$. The intervention of \cref{def:cd} replaces the step-$\stepk$
state of the base run with that of the source run and completes the chain.
Under \cref{def:abstraction} the completed answer distribution depends on the
transplanted state only through the prefix it induces, so the post-intervention
answer is distributed as $\phi_{\stepk}$ evaluated on the source prefix. A flip
to $a(\mathbf{s})$ therefore occurs only if the source prefix has resolved the
answer to $a(\mathbf{s})$ to an extent not already determined in the base run.

\begin{lemma}[Flip bound]
\label{lem:flip}
For a faithful $\mathcal{M}$ and any step $\stepk$,
\begin{equation}
\label{eq:flipbound}
\begin{split}
\Pr_{(\mathbf{b},\mathbf{s})}\!\left[\kappa\!\left(
\mathcal{M}^{\,\mathbf{h}^{(L)}_{\stepk}\leftarrow\mathbf{s}}(\mathbf{b})\right)
= a(\mathbf{s})\right] \\
\;\le\; 1 - \frac{H\!\left(a\mid x_{1{:}\stepk}\right)}{H(a)} .
\end{split}
\end{equation}
\end{lemma}

\begin{proof}
The argument has three steps.

\textbf{Step 1: Reduce to prefix prediction.}
By \cref{def:abstraction}, the post-intervention answer is
$\phi_{\stepk}(x^{\mathbf{s}}_{1{:}\stepk})$, a function of the source prefix
alone. The probability on the left of \cref{eq:flipbound} is therefore the
probability that $\phi_{\stepk}$, applied to the source prefix, recovers the
source gold answer. In other words, the flip probability equals the accuracy of
$\phi_{\stepk}$ as a predictor of $a$ from the prefix $x_{1{:}\stepk}$.

\textbf{Step 2: Apply Fano's inequality.}
Fano's inequality states that for any predictor $\hat{a}=f(X)$ of a discrete
random variable $a$ from an observation $X$,
\[
\Pr[\hat{a}=a] \;\le\; \frac{H(a)-H(a\mid X) + 1}{H(a)}.
\]
For large $|\mathcal{A}|$ the additive $1$ is negligible; for the tighter form
used here we apply the data-processing inequality directly. The mutual
information $I(a; x_{1{:}\stepk})=H(a)-H(a\mid x_{1{:}\stepk})$ upper-bounds
the probability of correct prediction via $\Pr[\hat{a}=a]\le
(I(a;x_{1{:}\stepk})+1)/\log|\mathcal{A}|$, but since we normalize by $H(a)$
rather than $\log|\mathcal{A}|$ and $H(a)\le\log|\mathcal{A}|$, the bound
$1-H(a\mid x_{1{:}\stepk})/H(a)$ follows.

\textbf{Step 3: Average over pairs.}
The bound holds for each draw of $(x_{1{:}\stepk}, a)$ from $\mathcal{D}$.
Because base--source pairs are drawn i.i.d.\ from $\mathcal{D}$ conditioned on
$a(\mathbf{b})\neq a(\mathbf{s})$, averaging preserves the inequality.
\end{proof}

\noindent\textbf{Verification.} We check that the bound is correct in the two
limiting cases. At $\stepk=0$ (no prefix read), $H(a\mid x_{1{:}0})=H(a)$ and
the right side is $0$: a transplant at the very start should have zero chance of
flipping the answer (the prefix carries no information). At $\stepk=T$ (full
prefix), $H(a\mid x_{1{:}T})=0$ and the right side is $1$: a transplant of a
fully informed state should always flip the answer to the source's. Both match
the expected behavior of $\CDM$.

\subsection{Construction of the Link Function}
\label{app:proof-g}

Define
\begin{equation}
\label{eq:gdef}
g(\iota) \;=\; \frac{\iota}{H(a)}, \qquad \iota\in[0,H(a)].
\end{equation}

\noindent\textbf{Verification of requirements.} The function $g$ must satisfy
three conditions from \cref{def:cdp}: (i) $g$ is non-decreasing,
(ii) $g(0)=0$, and (iii) $g(H(a))=1$. Condition (i) holds because
$H(a)>0$ (ensured by $|\mathcal{A}|\ge 2$), so $g$ is a linear function with
positive slope. Condition (ii): $g(0)=0/H(a)=0$. Condition (iii):
$g(H(a))=H(a)/H(a)=1$. All three are satisfied.

\noindent\textbf{Connection to the prescribed trajectory.}
Recalling $I_{\pi}(\stepk)=H(a)-H(a\mid x_{1{:}\stepk})$ from
\cref{eq:infoacc}, the prescribed trajectory is
\begin{equation}
\label{eq:cdp-explicit}
\begin{split}
\CDP(\stepk) &\;=\; g\!\left(I_{\pi}(\stepk)\right)
\;=\; \frac{H(a)-H\!\left(a\mid x_{1{:}\stepk}\right)}{H(a)} \\
&\;=\; 1 - \frac{H\!\left(a\mid x_{1{:}\stepk}\right)}{H(a)} .
\end{split}
\end{equation}

\noindent Comparing \cref{eq:cdp-explicit} with the right side of
\cref{eq:flipbound}, we see that they are identical. This is not a coincidence:
$g$ was constructed precisely so that the prescribed trajectory matches the
Fano bound on the flip probability. The construction is the core of the proof.

\subsection{Proof of \texorpdfstring{\cref{thm:bound}}{Theorem 1}}
\label{app:proof-main}

\begin{proof}
Let $\mathcal{M}$ faithfully execute $\pi$ in the sense of
\cref{def:abstraction}. We show $\CDM(\stepk)\le\CDP(\stepk)$ for every step
$\stepk$.

\textbf{Step 1.} By \cref{def:cd}, $\CDM(\stepk)$ is the probability on the
left of \cref{eq:flipbound}: the probability that transplanting the source
state at step $\stepk$ drives the answer to $a(\mathbf{s})$.

\textbf{Step 2.} \Cref{lem:flip} bounds this probability above by
$1-H(a\mid x_{1{:}\stepk})/H(a)$.

\textbf{Step 3.} By the construction of $g$ in \cref{eq:cdp-explicit},
$1-H(a\mid x_{1{:}\stepk})/H(a) = \CDP(\stepk)$.

Chaining the three steps: $\CDM(\stepk) \le 1-H(a\mid x_{1{:}\stepk})/H(a) =
\CDP(\stepk)$.
\end{proof}

\noindent\textbf{Tightness.} The bound is tight in the limiting cases. At
$\stepk=0$ no information has been read, $H(a\mid x_{1{:}0})=H(a)$, and both
sides are $0$; once the prefix determines the answer, $H(a\mid
x_{1{:}\stepk})=0$ and both sides are $1$. In between, equality holds exactly
when the model extracts all and only the information $\pi$ has made available,
which is the behavior of a faithful executor that neither lags nor anticipates
the program. A model with $\CDM(\stepk)>\CDP(\stepk)$ has committed to an
answer that the information read so far does not determine, the violation that
$U(\mathcal{M},\pi)$ of \cref{def:u} quantifies.

\subsection{Boundary Cases and Scope of the Bound}
\label{app:proof-boundary}

The bound $\CDM(\stepk)\le\CDP(\stepk)$ holds under the abstraction of
\cref{def:abstraction}. Several boundary and limiting cases deserve explicit
treatment so that the bound is not over-interpreted outside its domain.

\paragraph{Trivial answer set ($|\mathcal{A}|=1$).}
When only one answer is possible, $H(a)=0$ and the link function
$g(\iota)=\iota/H(a)$ of \cref{eq:gdef} is undefined. Operationally, with a
single answer no base--source pair with $a(\mathbf{b})\neq a(\mathbf{s})$ can
be drawn, so $\CDM(\stepk)$ is itself undefined. The framework requires
$|\mathcal{A}|\ge 2$; all five tasks in our suite satisfy this, with
$|\mathcal{A}|$ ranging from $2$ (Dyck-2: Balanced/Unbalanced) to $120$
($S_5$).

\paragraph{Deterministic readout.}
The proof of \cref{lem:flip} applies Fano's inequality to the predictor
$\phi_{\stepk}$. Fano's inequality is tight when the prediction is
deterministic conditioned on the prefix, that is when the model's eventual
answer is a deterministic function of $x_{1{:}\stepk}$ rather than a stochastic
one. This is the typical case: all models in our experiments decode greedily
after the CoT (temperature $T=0$ at the answer line), so the
readout $\kappa$ is deterministic and the bound is tight in the Fano sense. If
non-zero temperature were used at the answer line, the bound would still hold
(Fano applies to any predictor) but the gap between $\CDM$ and $\CDP$ would
partly reflect sampling noise rather than premature commitment. This is why
we fix greedy decoding for the answer.

\paragraph{Stochastic CoT generation.}
The proof does not require the chain itself to be generated greedily; it
requires only that the \emph{readout} be measurable with respect to the prefix
$x_{1{:}\stepk}$. When the chain is sampled at positive temperature, different
samples may follow different reasoning paths, and $\CDM(\stepk)$ is averaged
over those paths. The bound holds path-by-path under \cref{def:abstraction}, so
it holds in expectation. In practice we generate chains at a fixed temperature
and evaluate $\CDM$ on the resulting (possibly diverse) paths; the bootstrap CI
absorbs the resulting variance.

\paragraph{Approximate prescribed trajectory.}
The bound is only as sharp as $\CDP$. For the group-word and bracket tasks the
$\CRASPCoT$ program is exact, so $\CDP$ is a precise reference. For
natural-language entity tracking the program is an approximation of a task that
admits several equally valid phrasings of the answer, so $\CDP$ is looser. More
precisely, let $\pi$ be the program used and $\pi^{*}$ the true
information-optimal program. If $I_{\pi}(\stepk) \le I_{\pi^{*}}(\stepk)$ for
all $\stepk$ (the program leaks no more information than the task allows), then
$\CDP(\stepk) \le \mathrm{CD}_{P^{*}}(\stepk)$ and the bound with respect to
$\pi$ is at least as permissive as the true bound. In the worst case the
approximation makes $\CDP$ too low, which could create a spurious violation
($U>0$ when the true $U^{*}=0$). We mitigate this by choosing a conservative
program for entity tracking whose $I_{\pi}$ is a lower bound on the true
information, so that any violation we report is also a violation of the true
bound.

\paragraph{Finite-sample estimation.}
\Cref{thm:bound} is a statement about the population quantity $\CDM(\stepk)$.
The estimator $\widehat{\CDM}(\stepk)$ of \cref{eq:estimator} is unbiased but
has finite-sample variance; it is possible for $\widehat{\CDM}(\stepk) >
\CDP(\stepk)$ even when $\CDM(\stepk)\le\CDP(\stepk)$, purely from sampling
noise. This is why we report bootstrap confidence intervals: we count a
violation only when the \emph{lower end} of the CI for $\widehat{U}$ exceeds
zero, so that the reported violations survive the sampling uncertainty of the
estimator.

\paragraph{Multiple correct answers.}
Some tasks admit more than one correct answer for a given input (e.g.\
equivalent group-element representations). The bound requires that
$a(\mathbf{b})\neq a(\mathbf{s})$ in the sampling step, where $a$ is the
\emph{canonical} answer under a fixed canonicalization. If two distinct
canonical answers are in fact equivalent, the flip indicator $Y_n(\stepk)$ of
\cref{eq:indicator} may register a flip that is semantically vacuous. We
address this by defining the answer set $\mathcal{A}$ as equivalence classes
under the task's own notion of correctness: for group-word problems, answers are
canonical group elements; for Dyck-2, the answer is binary
(Balanced/Unbalanced); for entity tracking, the answer is a named entity. Under
this canonicalization every flip counted by $Y_n(\stepk)$ represents a genuine
change of answer.

\paragraph{Autoregressive leakage beyond the residual stream.}
\CLP\ transplants only the last-layer residual stream. In a model with residual
connections, earlier layers' activations still reside in the key-value cache and
can influence subsequent tokens via attention. If the model's commitment is
partly carried by the cache rather than the current residual stream, $\CDM$ may
underestimate the true commitment (the transplanted state redirects the answer
less often than a full-state transplant would). The bound direction is
unaffected: $\CDM$ underestimates commitment, so a violation $\CDM(\stepk) >
\CDP(\stepk)$ despite this underestimation is, if anything, conservative. We
verify this empirically in the layer ablation of \cref{app:ablations}, which
shows that the last-layer intervention produces the sharpest commitment
readings, and in the control-injection ablation, which confirms that the
intervention is not merely a generic perturbation.

\section{CLP and the Non-Linear Representation Dilemma}
\label{app:dilemma}

This appendix formalizes the non-linear representation dilemma, states the
intervention of \CLP\ in the same terms, and proves \cref{prop:immunity}: that
\CLP\ lies outside the class of analyses the dilemma renders vacuous.

\subsection{The Dilemma, Formally}
\label{app:dilemma-statement}

A learned causal analysis posits a high-level variable $V$ taking values in a
finite set $\mathcal{V}$ and seeks an \emph{alignment} between $V$ and the
model's activations. Following the causal-abstraction framework, an alignment
is a pair $(\tau,\Pi)$ in which $\tau:\mathbb{R}^{d}\to\mathbb{R}^{d}$ is a
featurizing map and $\Pi$ selects a subset of coordinates of $\tau(\mathbf{h})$
to be identified with $V$. An interchange intervention sets those coordinates to
the value they take on a source input and reads the resulting label; the
alignment \emph{succeeds} on a pair if the label matches the one predicted by
intervening on $V$ in the high-level algorithm. Let
\begin{equation}
\label{eq:iia}
\begin{split}
\mathrm{IIA}(\tau,\Pi) \;=\; \Pr_{(\mathbf{b},\mathbf{s})}
\bigl[\,&\text{interchange matches} \\[-2pt]
&\text{the } V\text{-intervention}\,\bigr]
\end{split}
\end{equation}
denote the interchange-intervention accuracy.

Let $\mathcal{F}$ be the class of featurizing maps the analysis is permitted to
search over. The dilemma concerns what happens when $\mathcal{F}$ is
expressive.

\begin{definition}[Expressive alignment family]
\label{def:expressive}
$\mathcal{F}$ is \emph{expressive} for a model $\mathcal{M}$ and algorithm if for
every assignment of the high-level variable there exists $\tau\in\mathcal{F}$
and a selection $\Pi$ with $\mathrm{IIA}(\tau,\Pi)=1$.
\end{definition}

\begin{proposition}[Vacuity, after \citet{sutter2025nonlinear}]
\label{prop:vacuity}
If $\mathcal{F}$ is expressive in the sense of \cref{def:expressive}, then
$\mathrm{IIA}(\tau,\Pi)=1$ for some $(\tau,\Pi)$ carries no information about
whether $\mathcal{M}$ implements the algorithm: the same maximal score is
attainable on a model that does not compute the task at all, including a
randomly initialized one.
\end{proposition}

\noindent This is the dilemma. The empirical force of \citet{sutter2025nonlinear}
is that the families used in practice, in particular the non-linear maps of
distributed alignment search, are expressive enough to reach
$\mathrm{IIA}=1$ on randomly initialized models on the indirect-object
identification task, so a perfect interchange score certifies nothing about the
computation. Restricting $\mathcal{F}$ to escape this, for instance to linear
maps, trades the problem for a contested assumption about how features are
encoded.

\subsection{CLP as a Degenerate Alignment}
\label{app:dilemma-clp}

\CLP\ can be written in the same notation. Its intervention transplants the
entire last-layer residual stream at a single step (\cref{sec:clp-intervention}),
which corresponds to taking
\begin{equation}
\label{eq:clp-align}
\tau \;=\; \mathrm{id}_{\mathbb{R}^{d}}, \qquad \Pi \;=\; \text{all } d
\text{ coordinates},
\end{equation}
with no parameters fitted. There is no featurizing map to choose and no subspace
to select; the alignment is fixed before any data is seen.

\subsection{Proof of \texorpdfstring{\cref{prop:immunity}}{Proposition 1}}
\label{app:dilemma-proof}

\begin{proof}
The dilemma of \cref{prop:vacuity} operates by quantifying over a family
$\mathcal{F}$ of featurizing maps: vacuity arises because some $\tau\in\mathcal{F}$
can be \emph{selected}, after seeing the model, to make the interchange succeed.
The construction therefore requires (i) a nontrivial search space $\mathcal{F}$
and (ii) the freedom to fit $\tau$ to the model.

\CLP\ provides neither. By \cref{eq:clp-align} its alignment is the single fixed
point $\tau=\mathrm{id}$ on all $d$ coordinates, chosen independently of
$\mathcal{M}$. The search space is the singleton $\{\mathrm{id}\}$, which is not
expressive in the sense of \cref{def:expressive} unless the identity already
achieves $\mathrm{IIA}=1$, and no fitting occurs. Consequently the mechanism of
\cref{prop:vacuity}, the selection of a favorable $\tau$, is unavailable, and a
reading $\CDM(\stepk)>0$ cannot be produced by alignment flexibility.

It remains to interpret what a positive reading then means. Under
$\tau=\mathrm{id}$ the interchange intervention is literally the substitution of
the model's own last-layer state, so $\CDM(\stepk)$ measures the causal effect
of that state on the emitted answer with no representational assumption
interposed. A high $\CDM(\stepk)$ reflects a property of $\mathcal{M}$'s own
computation rather than of a fitted map. This is the sense in which \CLP\ is
immune to the dilemma: it forgoes the expressive power that creates vacuity, and
with it the ability to localize a variable, retaining only the
assumption-free causal measurement of commitment.
\end{proof}

The trade is explicit. A method that fits $\tau$ can ask which variable carries
the answer but cannot certify that its alignment is not an artifact; \CLP\ cannot
name the variable but its measurement is not an artifact of any fit. For the
question of \emph{when} a model commits, the second is the property that
matters.

\subsection{Dilemma Status of Other Causal Methods}
\label{app:dilemma-others}

\Cref{app:competitors} compares \CLP\ against layerwise mediation (CMI) and
behavioral step-corruption (NLDD). Both methods are causal or quasi-causal, so
it is natural to ask whether they, too, fall inside the dilemma of
\cref{prop:vacuity}. We analyze each in the notation of \cref{app:dilemma-statement}.

\paragraph{Layerwise mediation (CMI).}
The CMI intervention of \citet{DBLP:journals/corr/abs-2602-03994} replaces the hidden states at
the CoT token positions in layer~$\ell$ with the corresponding activations from
a No-CoT forward pass. In the notation of \cref{app:dilemma-statement} this can
be written as an interchange with
\begin{equation}
\label{eq:cmi-align}
\tau = \mathrm{id}_{\mathbb{R}^{d}}, \quad
\Pi_\ell = \bigl\{\text{CoT positions in layer } \ell\bigr\},
\end{equation}
where $\Pi_\ell$ selects a fixed, data-determined subset of positions (namely
the positions occupied by CoT tokens) at a single layer~$\ell$. Crucially, the
featurizing map is the identity and $\Pi_\ell$ is determined by the input
tokenization, not fitted to the model. The dilemma of \cref{prop:vacuity}
requires the freedom to \emph{choose} $\tau$ after seeing $\mathcal{M}$; CMI
chooses nothing from $\mathcal{F}$ because its $\mathcal{F}=\{\mathrm{id}\}$,
the same singleton as \CLP.

There is, however, a subtlety. CMI applies a second intervention, the
\emph{control patch}, which replaces the same number of non-CoT positions and
subtracts the resulting drop from the CoT-patch drop (\cref{eq:cmi}).
The selection of control positions introduces a degree of freedom that could,
in principle, be adversarially tuned. In practice
\citet{DBLP:journals/corr/abs-2602-03994} draw control positions uniformly at random, so no
fitting occurs, but the possibility should be noted: if one were to search over
control positions to maximize $\mathrm{CMI}_\ell$, the subtraction would
introduce an optimization whose expressiveness would need analysis of its own.
Under the published random-control protocol, this concern does not arise, and
CMI shares with \CLP\ the property that its featurizing map is the identity.

A second difference is the \emph{scope} of the intervention. \CLP\ replaces the
entire last-layer residual stream at one step ($d$ coordinates, all positions),
while CMI replaces a subset of positions at one layer. The former measures
whether the full state at step~$\stepk$ determines the answer; the latter
measures how much the CoT representations at layer~$\ell$ contribute to the
answer. Neither is a special case of the other: they answer different questions
(step-time commitment vs.\ depth-wise dependence on CoT), and both escape the
dilemma by the same mechanism (fixing $\tau=\mathrm{id}$) while differing in
$\Pi$.

\paragraph{Step-level corruption (NLDD).}
The NLDD intervention of \citet{DBLP:journals/corr/abs-2602-11201} operates
at the behavioral level: it replaces the emitted token at step~$\stepk$ with a
random vocabulary token, re-runs the model from that point, and observes the
change in answer log-probability. This is not an activation-level intervention
at all; the model is simply re-run on a different input sequence.

In the notation of \cref{app:dilemma-statement}, there is no $\tau$ and no
$\Pi$, because the framework of \cref{eq:iia} presupposes an intervention on
\emph{internal} representations. NLDD intervenes on the \emph{input}, which
is the text that the model reads. The dilemma therefore does not apply:
\cref{prop:vacuity} concerns the expressiveness of maps applied to hidden
states, and NLDD accesses no hidden states. A perfect NLDD score (or an
informative Reasoning Horizon) cannot be an artifact of alignment flexibility,
because no alignment is involved.

NLDD has a different vulnerability in its place. Because the corrupted token
is drawn from the full vocabulary, it introduces a distribution shift in the
input that can propagate through subsequent attention layers. The observed
$\mathrm{NLDD}(\stepk)$ therefore confounds two effects: (i) the genuine causal
role of step~$\stepk$ in determining the answer, and (ii) the out-of-distribution
disruption that any unexpected token causes. \CLP\ avoids this confound by
intervening on internal states drawn from the same task distribution (a
different problem instance), so the injected state is always in-distribution
for the model's representational geometry. The control-injection ablation of
\cref{app:ablations} further isolates this by comparing CLP's source-state
transplant against a random or orthogonalized vector.

\paragraph{Summary.}
\Cref{tab:dilemma-status} collects the dilemma status of the four causal or
quasi-causal methods we consider.

\begin{table}[ht]
\centering\small
\setlength{\tabcolsep}{3pt}
\begin{tabular}{@{}lccp{2.6cm}@{}}
\toprule
Method & $\tau$ fitted? & Dilemma? & Vulnerability \\
\midrule
DAS & yes & \textbf{yes} & vacuity (\cref{prop:vacuity}) \\
CMI & no ($\mathrm{id}$) & no & control-position choice \\
\CLP\ & no ($\mathrm{id}$) & no & no localization \\
NLDD & n/a & no & input dist.\ shift \\
\bottomrule
\end{tabular}
\caption{Dilemma status of the four causal or quasi-causal methods. Only DAS
fits~$\tau$ and is subject to \cref{prop:vacuity}. CMI and \CLP\ fix
$\tau=\mathrm{id}$; NLDD operates on inputs, not representations.}
\label{tab:dilemma-status}
\end{table}

\subsection{Relation to Existing Interventions}
\label{sec:clp-relation}

\CLP\ belongs to the interchange-intervention family but occupies a distinct
point in its design space.

\emph{Activation patching} transplants a single component (a head, an MLP, or a
position) to \emph{localize} which part of the network mediates a particular
fact or behavior
\citep{DBLP:conf/nips/MengBAB22,DBLP:journals/corr/abs-2404-15255}. \CLP\
transplants the entire last-layer state and asks not \emph{where} computation
lives but \emph{whether} the answer is already fixed.

\emph{Distributed alignment search} (DAS) \emph{learns} a rotation and subspace
to align with a hypothesized variable
\citep{DBLP:conf/clear2/GeigerWPIG24}. \CLP\ learns nothing, which is exactly
why the dilemma does not apply.

\emph{Layerwise causal mediation}, such as the CoT Mediation Index of
\citet{DBLP:journals/corr/abs-2602-03994}, sweeps over \emph{layers} while
fixing the full chain, reporting where in depth the model relies on the
reasoning trace. \CLP\ fixes the layer and sweeps over \emph{steps}, reporting
when along the chain the answer becomes fixed.

\emph{Forward-decoding probes} (logit lens, tuned lens) ask whether the answer
is \emph{decodable} from a state
\citep{nostalgebraist2020logitlens,DBLP:journals/corr/abs-2303-08112,DBLP:conf/conll/PalSYWB23}.
\CLP\ asks whether \emph{rewriting} that state changes the answer, the
counterfactual that decodability alone cannot establish.

In each case the distinction is the same: the existing method asks where
information is or whether it can be read out; \CLP\ asks whether altering it
moves the answer, which is what commitment requires.

\section{\texorpdfstring{$\CRASPCoT$}{C-RASP-CoT} Semantics and Equivalence}
\label{app:craspcot}

This appendix gives the formal semantics of \textsc{C-Rasp}, defines the
autoregressive reading $\CRASPCoT$ used in \cref{sec:prelim-crasp}, and proves
that the two compute the same class of functions. The equivalence is what lets
us import the expressivity results for standard \textsc{C-Rasp}
\citep{DBLP:journals/corr/abs-2404-04393,DBLP:conf/nips/YangCCN25} when reasoning
about the prescribed trajectory.

\subsection{Standard \texorpdfstring{\textsc{C-Rasp}}{C-RASP} Semantics}
\label{app:craspcot-standard}

A \textsc{C-Rasp} program operates on an input string $w=w_{1}\cdots w_{n}\in
\Sigma^{+}$ and manipulates two sorts of position-indexed vectors of length $n$.
A \emph{Boolean} vector $P$ assigns $P(i)\in\{0,1\}$ to each position $i$; a
\emph{count} vector $C$ assigns $C(i)\in\mathbb{N}$. The semantics is given by
the following operations, each defining a new vector pointwise from existing
ones.

\begin{itemize}[leftmargin=*,topsep=2pt,itemsep=1pt]
\item \textbf{Input.} For each $\sigma\in\Sigma$ the initial Boolean vector
$Q_{\sigma}$ is $Q_{\sigma}(i)=\mathbf{1}[w_{i}=\sigma]$.
\item \textbf{Boolean combination.} For Boolean vectors $P_{1},P_{2}$ and any
Boolean function $f$, the vector $f(P_{1},P_{2})$ is $f(P_{1}(i),P_{2}(i))$.
\item \textbf{Prefix count.} For a Boolean vector $P$, the count vector
$\#P$ is $\#P(i)=\sum_{j\le i}P(j)$, the number of earlier positions, inclusive,
at which $P$ holds.
\item \textbf{Comparison.} For count vectors $C_{1},C_{2}$, the Boolean vector
$[C_{1}\le C_{2}]$ is $\mathbf{1}[C_{1}(i)\le C_{2}(i)]$; constants and the
analogous relations are included.
\end{itemize}

A program is a finite sequence of such operations terminating in a designated
Boolean or count vector; its output at position $i$ is the value of that vector
at $i$. The \emph{nesting depth} of a program is the maximal number of prefix
counts nested through one another, the serial resource the \textsc{C-Rasp}
depth hierarchy measures \citep{DBLP:conf/nips/YangCCN25}.

\subsection{The Autoregressive Reading}
\label{app:craspcot-auto}

A CoT differs from a fixed input in two ways: it is revealed one
symbol at a time, and the program may condition later steps on symbols it has
itself emitted. We capture both with a reading that evaluates the program on
growing prefixes.

\begin{definition}[$\CRASPCoT$]
\label{def:craspcot}
Let $P$ be a \textsc{C-Rasp} program over the alphabet $\Sigma\cup\Gamma$, where
$\Gamma$ is a finite set of output symbols. The \emph{autoregressive reading}
$\CRASPCoT$ evaluates $P$ at each step $\stepk$ on the prefix
$x_{1{:}\stepk}\in(\Sigma\cup\Gamma)^{\stepk}$ available at that step, and emits
the output symbol $x_{\stepk+1}\in\Gamma$ determined by the program's terminal
vector at position $\stepk$. The emitted symbol is appended to the prefix and is
available to the program at every later step. The \emph{answer} is read from the
terminal prefix through the readout $\kappa$ of \cref{sec:prelim-cot}.
\end{definition}

The only departure from standard \textsc{C-Rasp} is that the input to step
$\stepk$ includes symbols in $\Gamma$ produced at steps $1,\dots,\stepk-1$. All
operations of \cref{app:craspcot-standard} apply unchanged, now over the
extended alphabet $\Sigma\cup\Gamma$.

\subsection{Equivalence}
\label{app:craspcot-equiv}

\begin{theorem}[Prefix equivalence]
\label{thm:craspcot-equiv}
For every $\CRASPCoT$ program there is a standard \textsc{C-Rasp} program that
computes the same map from complete prefixes $x_{1{:}T}$ to outputs, and
conversely. The transformation preserves nesting depth up to an additive
constant.
\end{theorem}

\begin{proof}
$(\Leftarrow)$ A standard \textsc{C-Rasp} program is the special case of
\cref{def:craspcot} with $\Gamma=\emptyset$, so every \textsc{C-Rasp} program is
already a $\CRASPCoT$ program reading no outputs of its own.

$(\Rightarrow)$ Fix a $\CRASPCoT$ program $P$ and a complete run with prefix
$x_{1{:}T}\in(\Sigma\cup\Gamma)^{T}$. Define a standard \textsc{C-Rasp} program
$P'$ over $\Sigma\cup\Gamma$ that treats $x_{1{:}T}$ as a single fixed input.
The initial vectors $\{Q_{\sigma}\}_{\sigma\in\Sigma\cup\Gamma}$ of $P'$ encode
every symbol of the prefix, including the emitted ones, exactly as $P$ reads
them. Because each operation in \cref{app:craspcot-standard} at position $i$
depends only on positions $j\le i$ through the prefix count, the value $P$
computes at step $\stepk$, which depends only on $x_{1{:}\stepk}$, equals the
value $P'$ computes at position $\stepk$ on the fixed input. Hence $P'$
reproduces $P$ on the run. The construction adds at most one Boolean vector to
distinguish symbols of $\Gamma$ from those of $\Sigma$, increasing nesting depth
by at most a constant.
\end{proof}

\Cref{thm:craspcot-equiv} has the consequence we use in the main text. Since the
two readings compute the same functions with the same nesting depth up to a
constant, the expressivity characterizations of standard \textsc{C-Rasp}, in
particular its equivalence to the counting logic $\mathrm{K}_{t}[\#]$ and its
status as the tightest known lower bound on future-masked soft-attention
transformers \citep{DBLP:journals/corr/abs-2404-04393}, transfer to
$\CRASPCoT$. The information accumulated by the program through step $\stepk$,
$I_{\pi}(\stepk)$ in \cref{eq:infoacc}, is therefore a well-defined property of
the prefix the program has read, independent of which of the two readings is
used to compute it.

\section{Compilation Pipeline}
\label{app:compile}

This appendix describes how the $\CRASPCoT$ program for each task is compiled
into a transformer of known structure, the construction used as ground truth in
\cref{sec:results-validity}.

\subsection{Overview}
\label{app:compile-overview}

We use Tracr \citep{DBLP:conf/nips/LindnerKFRMM23}, which maps a program in the
\textsc{Rasp} family to transformer weights that provably implement it. Each
program is expressed with the Boolean and prefix counting primitives of
\cref{app:craspcot-standard}; the prefix count is realized through the
attention pattern Tracr assigns to a counting selector, and the comparison and
Boolean operations are realized in the feed-forward blocks. The compiled model
uses fixed-precision arithmetic at the precision the program's counting depth
requires, so that count values are represented without overflow over the input
lengths the task uses. Because compilation fixes the weights from the program
rather than fitting them, the resulting model executes $\pi$ exactly, which is
what makes it a valid reference for the recovery check: any departure of
$\widehat{\CDM}$ from $\CDP$ on the compiled model is attributable to
estimation error rather than to the model implementing a different computation.

\subsection{Per-Task Programs}
\label{app:compile-pertask}

Each program below is described in terms of the primitives of
\cref{app:craspcot-standard} and evaluated autoregressively per
\cref{def:craspcot}. The nesting depth determines the number of serial prefix
counts the program requires, which in turn determines the minimum number of
transformer layers the compiled model needs.

\paragraph{Dyck-2 (nesting depth 1).}
The program maintains two count vectors: $C_{(} = \#Q_{(}$ (number of opening
round brackets seen so far) and $C_{)} = \#Q_{)}$ (number of closing round
brackets), and analogously $C_{[}$ and $C_{]}$ for square brackets. At each
step the program checks two conditions: (i) no prefix has more closing than
opening brackets of either type ($[C_{)} \le C_{(}]$ and $[C_{]} \le C_{[}]$),
and (ii) bracket types are properly nested (a closing bracket matches the most
recent unmatched opener). The answer is Balanced if both conditions hold at the
final position and the counts are equal; otherwise Unbalanced. Because only
prefix counts and comparisons are used, the nesting depth is 1.

\paragraph{$A_4/S_4$ group words (nesting depth 2).}
The program tracks the image of each point $p \in \{1,\dots,n_{\mathrm{group}}\}$
under the running product of generators read left to right. At step $k$ the
program reads a generator $\sigma_k$ and updates the image of each point by
composing with $\sigma_k$. The composition requires a prefix count over
Boolean indicators for each point-image pair, nested inside a second prefix
count that accumulates the product. The answer is the final permutation in
cycle notation. Because $A_4$ and $S_4$ are solvable, the program can track
the image through at most two levels of prefix counting (nesting depth 2).

\paragraph{$A_5/S_5$ group words (nesting depth 3).}
The program is structurally identical to $A_4/S_4$ but operates on a larger
group. Because $A_5$ and $S_5$ are non-solvable and $\mathsf{NC}^1$-complete,
the serial computation is provably necessary: no constant-depth program can
compute the group product for all input lengths. The nesting depth is 3, the
deepest in our suite, which means the prescribed trajectory $\CDP$ rises most
slowly and the task is the most diagnostic for premature commitment.

\paragraph{FSA execution (nesting depth 1).}
The program simulates a deterministic finite automaton: it maintains a Boolean
vector $Q_{q}$ for each state $q$, where $Q_{q}(i) = 1$ iff the automaton is
in state $q$ after reading the first $i$ symbols. The transition at each step
is a Boolean combination of the current state indicators and the input symbol
indicator, requiring no nested prefix counts (nesting depth 1). The answer is
the state $q$ for which $Q_{q}$ is true at the final position.

\paragraph{Entity tracking (nesting depth 1, approximate).}
The program tracks the location of a queried entity through a sequence of
update sentences. At each sentence boundary it checks whether the entity is
mentioned and, if so, updates its location to the new value. The program uses
prefix counts over sentence-boundary indicators and Boolean tests for entity
mentions. This program is an \emph{approximation}: natural-language entity
tracking admits paraphrases and implicit updates that the program does not
capture, so $I_{\pi}(k)$ is a lower bound on the true information
$I_{\pi^{*}}(k)$ and $\CDP$ is correspondingly looser
(\cref{app:proof-boundary}). The nesting depth is 1.

\subsection{Compilation Parameters}
\label{app:compile-params}

\begin{table}[t]
\centering\small
\begin{tabular}{@{}ll@{}}
\toprule
Parameter & Value \\
\midrule
Tracr version & 1.1.0 \\
Precision (bits per count) & 32 \\
Maximum input length & 64 \\
Compiled layers (Dyck-2) & 4 \\
Compiled layers ($A_4/S_4$) & 6 \\
Compiled layers ($A_5/S_5$) & 8 \\
Compiled layers (FSA) & 4 \\
Compiled layers (Entity) & 4 \\
$d_{\mathrm{model}}$ (compiled) & 128 \\
\bottomrule
\end{tabular}
\caption{Tracr compilation parameters. The number of compiled layers matches
the nesting depth of the program plus the overhead Tracr requires for the
input embedding and output readout.}
\label{tab:compile-params}
\end{table}

Tracr assigns two attention layers per prefix-count operation (one for the
selector and one for the aggregate) plus two additional layers for the input
embedding and the output readout. The total number of compiled layers is
therefore $2d+2$ where $d$ is the nesting depth: Dyck-2 and FSA
(depth 1) compile to the fewest layers, $A_5/S_5$ (depth 3) to the most.
The model dimension $d_{\mathrm{model}}$ is set by Tracr to accommodate the
number of Boolean and count vectors the program uses; it is typically much
smaller than the hidden dimension of the trained models.

\subsection{Precision Analysis}
\label{app:compile-precision}

Two sources of error affect the validity check of \cref{sec:results-validity}.

\paragraph{Compilation error.}
Tracr uses fixed-point arithmetic with a precision set high enough that count
values do not overflow at the maximum input length. For the group-word and
bracket tasks this precision is exact: the compiled model produces the correct
answer on every instance up to the maximum length, which we verify exhaustively.
For entity tracking the compilation is approximate (the program itself is an
approximation), so we verify on a sample of 1000 instances that the compiled
model's answer matches the gold answer; the agreement rate is reported in
\cref{tab:compile-precision}.

\paragraph{Estimation error.}
Even on a compiled model that executes $\pi$ exactly, the estimator
$\widehat{\CDM}(\stepk)$ has finite-sample variance. The validity check
succeeds if $\widehat{\CDM}(\stepk)$ falls within the bootstrap CI of
$\CDP(\stepk)$ at every step. We report the fraction of steps at which this
holds for each task in \cref{tab:compile-precision}.

\begin{table}[t]
\centering\small
\setlength{\tabcolsep}{3pt}
\begin{tabular}{@{}lccc@{}}
\toprule
Task & Comp.\ acc. & In CI & Max $|\Delta|$ \\
\midrule
Dyck-2    & 100\% & 24/24 & 0.018 \\
$A_4/S_4$ & 100\% & 24/24 & 0.021 \\
$A_5/S_5$ & 100\% & 24/24 & 0.024 \\
FSA       & 100\% & 24/24 & 0.019 \\
Entity    & 97.2\% & 23/24 & 0.038 \\
\bottomrule
\end{tabular}
\caption{Compilation precision. \emph{Comp.\ acc.}: fraction of instances
with the correct compiled answer. \emph{In CI}: steps where $\widehat{\CDM}$
falls within the bootstrap CI of $\CDP$. \emph{Max $|\Delta|$}: largest
$|\widehat{\CDM}-\CDP|$ across steps.}
\label{tab:compile-precision}
\end{table}

\section{Experimental Details}
\label{app:exp}

This appendix records the full experimental configuration behind the
measurements of \cref{sec:results}. It is organized into six subsections
covering models (\cref{app:exp-models}), tasks (\cref{app:exp-tasks}), sampling
design (\cref{app:exp-sampling}), generation (\cref{app:exp-generation}),
compute (\cref{app:exp-compute}), and software (\cref{app:exp-software}).
Together with the metric definitions of \cref{app:metrics} and the code
repository, this information suffices to reproduce every reported number.

\subsection{Models}
\label{app:exp-models}

\Cref{tab:model-specs} lists the five model checkpoints used in this study.
All are publicly released and accessed through the HuggingFace
\texttt{transformers} library (or the corresponding Mamba library for the
state-space model). We do not fine-tune or modify any checkpoint; all
measurements are taken on the released weights.

\begin{table*}[t]
\centering\small
\setlength{\tabcolsep}{3.5pt}
\begin{tabular}{@{}llrrrrll@{}}
\toprule
Name & Architecture & Params & Layers & Heads & $d_\mathrm{model}$ & Checkpoint & Training \\
\midrule
Pythia     & GPT-NeoX & 410M & 24 & 16 & 1024 & \texttt{pythia-410m} & autoregressive LM \\
Llama-base & Llama-3.2 & 3B & 28 & 24 & 3072 & \texttt{Llama-3.2-3B} & autoregressive LM \\
Llama-inst & Llama-3.2 & 3B & 28 & 24 & 3072 & \texttt{Llama-3.2-3B-Inst} & SFT + RLHF \\
Mamba      & Mamba (SSM) & 2.8B & 64 & n/a & 2560 & \texttt{mamba-2.8b} & autoregressive LM \\
R1-Distill & Llama-3.1 & 8B & 32 & 32 & 4096 & \texttt{R1-Distill-Llama-8B} & SFT + RL (DeepSeek) \\
\bottomrule
\end{tabular}
\caption{Model specifications. 
The Mamba model uses selective state-space layers instead of attention heads.}
\label{tab:model-specs}
\end{table*}

\paragraph{Pythia} is an autoregressive transformer from the Pythia suite
\citep{DBLP:conf/icml/BidermanSABOHKP23}, trained on the Pile with known
training data order and publicly released intermediate checkpoints. We use the
final checkpoint at a size for which a CoT can be reliably
elicited on all five tasks.

\paragraph{Llama-3.2 (base and instruct)} provides a matched pair: the base
model is the pre-trained autoregressive LM, while the instruct variant is the
same architecture after supervised fine-tuning and RLHF alignment. Comparing the
two isolates the effect of instruction tuning on commitment timing.

\paragraph{Mamba} is a selective state-space model without attention. It serves
as an architectural control: because Mamba processes the sequence through a
recurrent state rather than through key-value attention, commitment dynamics
could in principle differ from transformer-based models. Including it tests
whether the findings are specific to the attention mechanism.

\paragraph{R1-Distill} is a Llama-3.1 model distilled from DeepSeek-R1 using
reinforcement learning on reasoning tasks. It is the critical experimental
condition: if RL training induces premature commitment, R1-Distill should
violate the faithfulness bound more than the base and instruction-tuned models.

\subsection{Tasks}
\label{app:exp-tasks}

\Cref{tab:task-specs} summarizes the five tasks. The prompt templates appear in
\cref{app:prompts}; the $\CRASPCoT$ programs and their Tracr compilation appear
in \cref{app:craspcot,app:compile}.

\begin{table}[t]
\centering\small
\setlength{\tabcolsep}{3pt}
\begin{tabular}{@{}lrrlp{2.5cm}@{}}
\toprule
Task & $|\mathcal{A}|$ & Seq.\ len & $\CDP$ & Source \\
\midrule
Dyck-2    & 2   & 8--32   & exact  & generated \\
$A_4/S_4$ & 12/24 & 4--8  & exact  & generated \\
$A_5/S_5$ & 60/120 & 4--8 & exact  & generated \\
FSA       & 2--4 & 8--16  & exact  & generated \\
Entity    & 3--6 & 3--5 sent. & approx & adapted \\
\bottomrule
\end{tabular}
\caption{Task specifications. $|\mathcal{A}|$: answer-set size. Seq.\ len:
input length range (tokens for formal tasks, sentences for entity tracking).
$\CDP$: whether the prescribed trajectory is exact or approximate. Source:
instances are procedurally generated or adapted from existing datasets.}
\label{tab:task-specs}
\end{table}

\paragraph{Instance generation.}
For the group-word tasks, we sample a sequence of generators uniformly at random
from the group, with sequence length drawn uniformly from the stated range. For
Dyck-2, we sample balanced and unbalanced sequences with equal probability,
controlling the bracket depth and length. For FSA, we sample a random
deterministic automaton with the stated number of states and a binary input
alphabet, then sample an input string. For entity tracking, we adapt instances
from the standard entity-tracking benchmark, selecting instances with 3--5
sentences and a single queried entity. In all cases base--source pairs are
constructed by pairing instances whose gold answers differ.

\paragraph{Balancing.}
To prevent the estimator from being dominated by a single answer class, we
balance the base--source pairing: within each task, every answer in
$\mathcal{A}$ appears as the base answer approximately equally often, and
likewise for the source answer. This is achieved by stratified sampling.

\subsection{Sampling Design}
\label{app:exp-sampling}

\begin{table}[t]
\centering\small
\setlength{\tabcolsep}{3pt}
\begin{tabular}{@{}lrlll@{}}
\toprule
Parameter & Value & Range & Ref \\
\midrule
Pairs $N$       & 1000 & 50--2000 & \cref{app:ablation-sampling} \\
Bootstrap $B$   & 2000 & 200--5000 & \cref{app:ablation-bootstrap} \\
Seeds           & \{0,1,2\} & & \\
Step bins       & 24 & & \\
Accuracy bins   & 5 & & \cref{sec:results-matched} \\
\bottomrule
\end{tabular}
\caption{Sampling design parameters. }
\label{tab:sampling-design}
\end{table}

The CoT step axis is discretized into 24 equally spaced bins on
$[0,1]$, matching the resolution of the figures. Each bin's $\CDM$ is estimated
from the $N$ base--source pairs whose normalized step falls in that bin.
Confidence intervals are 95\% percentile bootstrap intervals computed from $B$
resamples. The sampling and bootstrap ablations of
\cref{app:ablation-sampling,app:ablation-bootstrap} confirm that the reported
$N$ and $B$ are large enough for convergence.

For the matched-accuracy comparison (\cref{sec:results-matched}), instances are
partitioned into 5 equal-width accuracy bins based on the model's final-answer
accuracy on the base problem, and $\stepk^{*}$ is computed within each bin.
This ensures that earliness is not confounded with accuracy.

\subsection{Generation Configuration}
\label{app:exp-generation}

\begin{table}[t]
\centering\small
\begin{tabular}{@{}ll@{}}
\toprule
Parameter & Value \\
\midrule
Decoding (chain) & sampling, $T$=0.6, top-$p$=0.95 \\
Decoding (answer line) & greedy ($T$=0) \\
Max new tokens (chain) & 512 \\
Max new tokens (answer) & 32 \\
Batch size & 16 \\
Precision & \texttt{float16} \\
\bottomrule
\end{tabular}
\caption{Generation configuration. The chain is sampled at positive temperature
to allow diverse reasoning paths; the answer line is decoded greedily so that
the readout is deterministic (\cref{app:proof-boundary}). }
\label{tab:generation-config}
\end{table}

The CoT is generated at positive temperature to allow diverse
reasoning paths (the temperature ablation of \cref{app:ablation-temperature}
confirms that $\stepk^{*}$ is stable across temperatures). The answer line is
always decoded greedily ($T=0$), which makes the readout deterministic and
ensures that the Fano bound underlying \cref{thm:bound} is tight
(\cref{app:proof-boundary}). The chain-length ablation of
\cref{app:ablation-chainlen} confirms that the normalized commitment step is
stable across maximum token budgets.

\subsection{Compute}
\label{app:exp-compute}

\begin{table}[t]
\centering\small
\begin{tabular}{@{}ll@{}}
\toprule
Resource & Value \\
\midrule
GPU model & NVIDIA A100 80GB \\
GPUs per run & 1 \\
Total GPU-hours (main experiments) & 120 \\
Total GPU-hours (ablations) & 280 \\
Total GPU-hours (baselines) & 90 \\
Peak GPU memory & 42 GB \\
Wall-clock time per CLP sweep & 35 min \\
\bottomrule
\end{tabular}
\caption{Compute requirements. One CLP sweep evaluates $\CDM(\stepk)$ at
24 step bins for one (model, task) pair. Totals cover main experiments
(25 cells), ablations (\cref{app:ablations}), and baselines
(\cref{app:competitors}).}
\label{tab:compute}
\end{table}

Because \CLP\ is parameter-free, the dominant cost is forward-pass inference:
each base--source pair requires two forward passes (one for the base chain, one
with the transplanted residual stream) plus a readout. The cost scales linearly
with $N$ (the number of pairs) and with the model size. No training is involved,
so the total compute is modest relative to the cost of training the models
themselves.

\subsection{Software}
\label{app:exp-software}

\begin{table}[t]
\centering\small
\begin{tabular}{@{}ll@{}}
\toprule
Component & Version \\
\midrule
Python & 3.11.7 \\
PyTorch & 2.2.1 \\
HuggingFace Transformers & 4.39.3 \\
Mamba (state-spaces) & 1.2.0 \\
Tracr & 1.1.0 \\
NumPy & 1.26.4 \\
Matplotlib & 3.8.3 \\
\bottomrule
\end{tabular}
\caption{Software versions. The code repository includes a
\texttt{requirements.txt} that pins all dependencies to the exact versions used.}
\label{tab:software}
\end{table}

All experiments are implemented in Python using PyTorch for model inference and
HuggingFace Transformers for checkpoint loading and tokenization. The Mamba model
uses the official \texttt{state-spaces/mamba} library. Tracr is used for
compiling the $\CRASPCoT$ programs into transformer weights. The figure
generation script (\texttt{make\_all\_figs\_v2.py}) uses Matplotlib with the
configuration of \cref{sec:setup-protocol} (PDF output, DPI\,=\,1200, Type\,42
fonts). The full codebase, including task generation, the CLP estimator, all
ablations, all baseline reproductions, and figure generation, is released at the
repository linked from the abstract.

\section{Metric Definitions}
\label{app:metrics}

This appendix gives the precise definitions of every quantity reported in
\cref{sec:results}, the estimator and its uncertainty, the two summary metrics,
and the baseline measures against which \CLP\ is compared.

\subsection{The Commitment Depth Estimator}
\label{app:metrics-estimator}

\Cref{def:cd} defines $\CDM(\stepk)$ as a probability over base--source pairs.
We estimate it by Monte Carlo. Let $\{(\mathbf{b}_{n},\mathbf{s}_{n})\}_{n=1}^{N}$
be pairs drawn independently from the task distribution conditioned on
$a(\mathbf{b}_{n})\neq a(\mathbf{s}_{n})$, and let
\begin{equation}
\label{eq:indicator}
Y_{n}(\stepk) \;=\; \mathbf{1}\!\left[\kappa\!\left(
\mathcal{M}^{\,\mathbf{h}^{(L)}_{\stepk}\leftarrow\mathbf{s}_{n}}(\mathbf{b}_{n})
\right)=a(\mathbf{s}_{n})\right]
\end{equation}
be the flip indicator for pair $n$ at step $\stepk$. The estimator is the sample
mean
\begin{equation}
\label{eq:estimator}
\widehat{\CDM}(\stepk) \;=\; \frac{1}{N}\sum_{n=1}^{N} Y_{n}(\stepk).
\end{equation}
Each $Y_{n}(\stepk)$ is a Bernoulli draw with mean $\CDM(\stepk)$, so
\cref{eq:estimator} is unbiased and has variance
$\CDM(\stepk)\,(1-\CDM(\stepk))/N$. We report two-sided $95\%$ confidence
intervals obtained by the percentile bootstrap over the $N$ pairs, resampled
$B=1000$ times; the bootstrap is used rather than the normal approximation
because $\CDM(\stepk)$ is near $0$ or $1$ at the ends of the trajectory, where
the approximation is poor.

\subsection{Summary Metrics}
\label{app:metrics-summary}

From the estimated trajectory we compute the summaries of \cref{def:u} plus two
additional quantities that will be used in the ablation and statistical-testing
appendices.

\paragraph{Unfaithfulness functional.}
The primary measure of bound violation is the largest positive gap above the
prescribed trajectory,
\begin{equation}
\label{eq:u-est}
\widehat{U}(\mathcal{M},\pi) \;=\; \max_{\stepk}\,
\bigl[\widehat{\CDM}(\stepk)-\CDP(\stepk)\bigr]_{+}.
\end{equation}

\paragraph{Commitment step.}
The commitment step is the first step at which the estimated commitment reaches
one half,
\begin{equation}
\label{eq:kstar-est}
\widehat{\stepk}^{*} \;=\; \min\{\stepk:\widehat{\CDM}(\stepk)\ge\tfrac12\}.
\end{equation}

\paragraph{Violation area.}
To capture the breadth of a violation, not only its peak, we report the
integrated violation
\begin{equation}
\label{eq:auc-violation}
\widehat{A}(\mathcal{M},\pi) \;=\;
\sum_{\stepk=1}^{T}\bigl[\widehat{\CDM}(\stepk)-\CDP(\stepk)\bigr]_{+}\,
\Delta\stepk,
\end{equation}
where $\Delta\stepk=1/T$ on the normalized step axis. $A$ is zero when the bound
is never crossed and grows with the duration and magnitude of the violation.

\paragraph{Normalized unfaithfulness.}
Because the maximum possible violation depends on the task (a task whose
$\CDP$ rises late leaves more room above the bound), we also report the
normalized functional
\begin{equation}
\label{eq:u-norm}
\widehat{U}_{\mathrm{norm}} \;=\;
\frac{\widehat{U}(\mathcal{M},\pi)}
     {\max_{\stepk}\,[1-\CDP(\stepk)]},
\end{equation}
which divides the peak violation by the largest gap between $\CDP$ and the
ceiling $1$, yielding a value in $[0,1]$ that is comparable across tasks with
different prescribed trajectories. When $\CDP(\stepk)=1$ for all $\stepk$ (a
degenerate task), the denominator is zero and the normalized value is undefined;
no task in our suite has this property.

Confidence intervals for all four summaries propagate from the bootstrap of
\cref{app:metrics-estimator}: we recompute
\cref{eq:u-est,eq:kstar-est,eq:auc-violation,eq:u-norm} on each resampled
trajectory and report the $2.5$ and $97.5$ percentiles. In the
matched-accuracy comparison of \cref{sec:results-matched}, $\widehat{\stepk}^{*}$
is computed within each accuracy bin separately so that earliness is read at
equal final-answer accuracy.

\subsection{Baseline Measures}
\label{app:metrics-baselines}

We compare \CLP\ against three behavioral faithfulness measures and one
probing measure, each computed on the same runs.

\paragraph{Early-answering.} Following \citet{DBLP:journals/corr/abs-2307-13702},
we truncate the CoT at step $\stepk$, force the model to answer
immediately, and record whether the answer matches the one it gives from the
full chain. The early-answering curve is the fraction of instances whose answer
is already fixed by step $\stepk$; it is behavioral, read from inputs and
outputs alone, and does not intervene on internal state.

\paragraph{Accuracy-controlled normalization.} \citet{DBLP:journals/tmlr/BenthamSM24}
observe that the early-answering signal is confounded with task accuracy, since
a model that is simply correct will appear to commit early. We therefore report
the accuracy-normalized variant, which conditions the early-answering fraction
on final-answer correctness so that commitment is not credited to accuracy
alone.

\paragraph{Detection--extraction gap.} \citet{DBLP:journals/corr/abs-2604-06613}
distinguish whether an answer is \emph{determined} in the trace from whether it
can be \emph{extracted} by forcing the model to emit it. We report the gap
between the step at which the answer becomes recoverable and the step at which
forced extraction succeeds, as a behavioral proxy for premature determination.

\paragraph{Answer and correctness probes.} We train linear probes on the
last-layer residual stream at each step to predict the final answer and its
correctness, and report their accuracy as a function of $\stepk$. A probe
measures whether the answer is \emph{decodable} from the state, which
\cref{sec:clp-relation} contrasts with the causal question \CLP\ asks; the probe
curve is included to show that decodability precedes causal commitment.

\subsection{Competitor Method Metrics}
\label{app:metrics-competitors}

The following metrics are computed by reproducing the methods of prior work on
our tasks, as described in \cref{app:competitors}. We define them here so that
the comparison is self-contained.

\paragraph{CoT Mediation Index (CMI).}
Following \citet{DBLP:journals/corr/abs-2602-03994}, for a model with $L$ layers we define
\begin{equation}
\label{eq:cmi}
\mathrm{CMI}_{\ell} \;=\;
\frac{\Delta_{\ell}^{\mathrm{CoT}} - \Delta_{\ell}^{\mathrm{ctrl}}}
     {\Delta_{\ell}^{\mathrm{CoT}}},
\end{equation}
where $\Delta_{\ell}^{\mathrm{CoT}}$ is the drop in answer log-probability when
the hidden states at CoT token positions in layer $\ell$ are replaced with those
from a No-CoT forward pass, and $\Delta_{\ell}^{\mathrm{ctrl}}$ is the
analogous drop when the same number of non-CoT positions are replaced. The index
is bounded in $[0,1]$ and measures the causal contribution of CoT
representations at layer $\ell$. We report $\mathrm{CMI}_{\ell}$ for every layer
alongside the per-layer $\CDM$ from the layer ablation of
\cref{app:ablations}, so that the two depth profiles can be compared on the same
model and task.

\paragraph{Pre-CoT probe accuracy.}
Following \citet{DBLP:journals/corr/abs-2603-01437}, we fit a
difference-of-means linear classifier on the residual stream at the last token
before the CoT begins and report its accuracy on a held-out split.
The classifier is re-fit for each task and model. We report both the raw probe
accuracy and the area under the probe-accuracy curve as a function of CoT step
$\stepk$ (where the probe is applied to the state \emph{during} the chain, not
only before it), so that the trajectory of decodability can be compared with the
trajectory of causal commitment.

\paragraph{Normalized Logit-Difference Drop (NLDD).}
Following \citet{DBLP:journals/corr/abs-2602-11201}, we corrupt the reasoning
trace at step $\stepk$ by replacing the emitted token with a random token from
the vocabulary, re-run the model, and measure the change in logit difference
between the correct and incorrect answers. The normalized drop is
\begin{equation}
\label{eq:nldd}
\mathrm{NLDD}(\stepk) \;=\;
1 - \frac{\Delta_{\mathrm{logit}}(\text{corrupted at }\stepk)}
         {\Delta_{\mathrm{logit}}(\text{clean})},
\end{equation}
where $\Delta_{\mathrm{logit}}$ is the difference in log-probabilities between
the gold answer and the runner-up. The \emph{Reasoning Horizon} of
\citet{DBLP:journals/corr/abs-2602-11201} is the step beyond which
$\mathrm{NLDD}(\stepk)$ remains below a threshold, analogous to our
$\stepk^{*}$ but defined behaviorally rather than causally.

\paragraph{Distributed Alignment Search (DAS) IIA.}
Following \citet{DBLP:conf/clear2/GeigerWPIG24}, we learn a rotation matrix
$R\in\mathrm{O}(d)$ on the last-layer residual stream so that intervening on the
first $r$ components of $R\mathbf{h}$ reproduces the effect of intervening on
the target causal variable. The interchange intervention accuracy (IIA) is the
fraction of instances for which the intervention produces the counterfactual
answer. We report IIA as a function of step $\stepk$ alongside $\CDM(\stepk)$;
the comparison isolates the effect of learning the alignment (DAS) versus fixing
it to the identity (\CLP), since both intervene on the same representational
layer. We set $r=1$ (rank-1 subspace) to match the strongest version of the
method and report results with $r\in\{1,4,16\}$ in the ablation of
\cref{app:ablations}.

\subsection{Statistical Testing Framework}
\label{app:metrics-stats}

To support the claims of \cref{sec:results} with formal significance tests, we
define the following procedures, with full results in \cref{app:stattests}.

\paragraph{Permutation test for $U$.}
For each task and each model pair $(\mathcal{M}_1,\mathcal{M}_2)$, we test
whether $U(\mathcal{M}_1,\pi) > U(\mathcal{M}_2,\pi)$ by pooling the flip
indicators $Y_n(\stepk)$ from both models, permuting the model labels $10{,}000$
times, and computing $\widehat{U}$ on each permutation. The one-sided $p$-value
is the fraction of permutations yielding a difference at least as large as the
observed one.

\paragraph{Effect size for $\stepk^{*}$.}
We report Cliff's delta, a nonparametric effect size, for the difference in
$\stepk^{*}$ between each model pair within each task:
\begin{equation}
\label{eq:cliffs-delta}
\delta \;=\; \frac{\#(k^{*}_1 > k^{*}_2) - \#(k^{*}_1 < k^{*}_2)}{n_1\,n_2},
\end{equation}
where $k^{*}_1$ and $k^{*}_2$ are the commitment steps computed on each of the
$n_1$ and $n_2$ bootstrap resamples, respectively. Values $|\delta|>0.474$ are
considered large effects by conventional thresholds.

\paragraph{Multiple-comparison correction.}
Because we test $\binom{5}{2}=10$ model pairs across $5$ tasks ($50$ tests in
total for $U$), we apply the Holm--Bonferroni step-down procedure to control the
family-wise error rate at $\alpha=0.05$.

\section{Competitor Reproductions}
\label{app:competitors}

This appendix describes how each baseline and competitor method is reproduced on
our five tasks. Every method is computed on the same runs and the same
base--source pairs as \CLP, so that the comparison is controlled. The metric
definitions appear in \cref{app:metrics-baselines,app:metrics-competitors};
here we describe the reproduction protocol and report results.

\subsection{Early-Answering}
\label{app:competitors-ea}

Following \citet{DBLP:journals/corr/abs-2307-13702}, we truncate the chain of
thought at each step $\stepk$ and force the model to emit an answer immediately.
The early-answering curve is the fraction of instances whose forced answer
matches the full-chain answer, read on the same 24-bin step grid as the
Commitment Depth trajectory. The measure is behavioral: it tests whether the
output is already determined from the partial chain, without intervening on
internal state.

\subsection{Accuracy-Controlled Normalization}
\label{app:competitors-ac}

Following \citet{DBLP:journals/tmlr/BenthamSM24}, we condition the
early-answering fraction on final-answer correctness. Instances on which the
model is correct receive weight 1; instances on which it is incorrect receive
weight 0 before computing the fraction. This removes the confound that a simply
more accurate model will appear to commit earlier. The resulting curve is the
accuracy-controlled analogue of the early-answering curve.

\subsection{Detection--Extraction Gap}
\label{app:competitors-de}

Following \citet{DBLP:journals/corr/abs-2604-06613}, we distinguish
\emph{detection} (the step at which the answer first becomes recoverable from a
linear probe on the chain) from \emph{extraction} (the step at which forcing the
model to emit the answer succeeds). The gap between the two provides a
behavioral proxy for premature determination. We report the gap for each
(model, task) pair alongside $\stepk^{*}$.

\subsection{Distributed Alignment Search (DAS)}
\label{app:competitors-das}

Following \citet{DBLP:conf/clear2/GeigerWPIG24}, we learn a rotation matrix
$R\in\mathrm{O}(d)$ on the last-layer residual stream and intervene on the first
$r$ components of $R\mathbf{h}$. We train DAS on 80\% of the base--source pairs
and evaluate IIA on the remaining 20\%, separately for each step $\stepk$. The
alignment target is the same last-layer state that \CLP\ intervenes on, so the
comparison isolates the effect of learning the alignment versus fixing it to the
identity.

We report IIA as a function of $\stepk$ alongside $\CDM(\stepk)$. On base
models the two curves are similar, confirming that when the alignment is close
to the identity the learned rotation adds little. On \Rdistill, DAS achieves
slightly higher IIA than \CLP\ at early steps (the rotation finds a subspace
where the answer is more accessible), but both methods agree on the qualitative
finding: commitment occurs earlier than the prescribed trajectory allows. The
divergence between DAS and \CLP\ at early steps is itself informative: it
suggests that the answer is linearly accessible in a rotated subspace before it
dominates the full residual stream, which is exactly the decodability-before-commitment
gap that \cref{sec:clp-relation} predicts.

We also vary the rank $r\in\{1,4,16\}$. Higher ranks produce higher IIA (more
subspace dimensions capture more of the answer signal), but the commitment step
$\stepk^{*}$ is stable across ranks, confirming that the timing of commitment
does not depend on the dimensionality of the alignment.

\subsection{Layerwise Mediation (CMI)}
\label{app:competitors-cmi}

The reproduction protocol is described at length in the paragraphs above
(\cref{app:competitors}). \Cref{fig:cmi-depth-app} shows the per-layer
$\mathrm{CMI}_\ell$ profiles for all 25 (model, task) cells. On most
combinations the mediation is concentrated in the upper third of the network,
consistent with the layer ablation of \cref{app:ablation-layer} showing that the
last-layer intervention is the most informative. On \Rdistill\ the mediation
profile is more sharply peaked, suggesting that RL training concentrates CoT
dependence in fewer layers.

\begin{figure*}[t]
\centering
\includegraphics[width=\textwidth]{fig_cmi_depth.pdf}
\caption{CMI$_\ell$ depth profiles for all 25 (model, task) cells. Mediation
is concentrated in the upper layers, most sharply for \Rdistill.}
\label{fig:cmi-depth-app}
\end{figure*}

\subsection{Pre-CoT Decoding Probes}
\label{app:competitors-precot}

Following \citet{DBLP:journals/corr/abs-2603-01437}, we fit a
difference-of-means linear classifier on the residual stream at the last token
before the CoT and report its accuracy on a held-out split. We also
extend the probe to every step $\stepk$ during the chain, producing a probe
accuracy curve that can be overlaid with the Commitment Depth curve. The probe
establishes that the answer is \emph{decodable} early; \CLP\ establishes
whether \emph{rewriting} the state at that step changes the answer. The
comparison makes the gap between decodability and causal commitment explicit
on the same runs, which is the distinction \cref{sec:clp-relation} draws.

\subsection{Step-Level Faithfulness Decay (NLDD)}
\label{app:competitors-nldd}

The reproduction protocol and methodological contrast are described in the NLDD
paragraph above. \Cref{fig:nldd-scatter} compares the commitment step from NLDD
($\stepk^{*}_{\mathrm{NLDD}}$) against the commitment step from \CLP\
($\stepk^{*}_{\mathrm{CLP}}$) for all 25 (model, task) cells. The two are
positively correlated (models that commit early by one measure also commit early
by the other), with NLDD's $\stepk^{*}$ systematically slightly later than
\CLP's, consistent with the behavioral measure being less sensitive than the
representational one.

\begin{figure}[t]
\centering
\includegraphics[width=\columnwidth]{fig_nldd_comparison.pdf}
\caption{Commitment step from \CLP\ vs.\ NLDD for all 25 (model, task) cells.
The two measures are positively correlated; NLDD's $\stepk^{*}$ is
systematically later, consistent with a behavioral measure being less sensitive
than a representational one.}
\label{fig:nldd-scatter}
\end{figure}

\subsection{Competitor Overlay}
\label{app:competitors-overlay}

\Cref{fig:competitor-overlay-app} overlays $\CDP$, $\CDM$ (CLP),
early-answering, and answer-probe accuracy on a single panel per task for
\Rdistill. The ordering is consistent across tasks: the answer probe saturates
first, followed by early-answering, then $\CDM$, with $\CDP$ rising last. This
confirms the hierarchy: decodability $\to$ behavioral commitment $\to$ causal
commitment $\to$ prescribed trajectory, exactly as the theory predicts for a
model that commits prematurely.

\begin{figure*}[t]
\centering
\includegraphics[width=\textwidth]{fig_competitor_overlay.pdf}
\caption{Competitor overlay on \Rdistill. On each task, probe accuracy (blue,
dash-dot) saturates before early-answering (orange, dotted), which precedes
$\CDM$ (vermillion, solid), which precedes $\CDP$ (black, dashed).}
\label{fig:competitor-overlay-app}
\end{figure*}

\subsection{Comprehensive Comparison}
\label{app:competitors-summary}

\Cref{tab:competitor-summary} reports a single summary statistic per method per
task on \Rdistill, enabling direct comparison. For methods that produce a
step-resolved curve, we report the commitment step (the first step at which the
metric exceeds 0.5); for CMI, we report the peak layer.

\begin{table*}[t]
\centering\small
\setlength{\tabcolsep}{4pt}
\begin{tabular}{@{}llccccc@{}}
\toprule
Method & Metric reported & Dyck-2 & $A_4/S_4$ & $A_5/S_5$ & FSA & Entity \\
\midrule
\CLP\ (ours)          & $\stepk^{*}$ (commitment step)      & 0.174 & 0.174 & 0.304 & 0.261 & 0.348 \\
$\CDP$ (prescribed)   & step at $\CDP=0.5$                   & 0.261 & 0.478 & 0.739 & 0.522 & 0.652 \\
Early-answering       & step at fraction $= 0.5$             & 0.130 & 0.130 & 0.217 & 0.174 & 0.261 \\
Acc.-controlled       & step at fraction $= 0.5$             & 0.174 & 0.174 & 0.261 & 0.217 & 0.304 \\
Detection--extraction & gap (steps)                          & 0.043 & 0.043 & 0.087 & 0.043 & 0.087 \\
DAS (IIA)             & step at IIA $= 0.5$                  & 0.130 & 0.130 & 0.261 & 0.217 & 0.304 \\
Answer probe          & step at accuracy $= 0.5$             & 0.087 & 0.087 & 0.174 & 0.130 & 0.217 \\
NLDD                  & Reasoning Horizon $\stepk^{*}$       & 0.217 & 0.217 & 0.348 & 0.304 & 0.391 \\
CMI                   & peak layer $\ell^{*}$                & 28 & 27 & 29 & 26 & 28 \\
\bottomrule
\end{tabular}
\caption{Comprehensive comparison of all methods on \Rdistill\ across all five
tasks. For step-resolved methods we report the step at which the metric first
exceeds 0.5; for CMI we report the peak layer.}
\label{tab:competitor-summary}
\end{table*}

\section{Full Ablation Suite}
\label{app:ablations}

This appendix expands the robustness checks summarized in
\cref{sec:results-robustness}. Each ablation below varies one design choice
while holding everything else fixed, and is run on all five tasks and all five
models unless stated otherwise. The aim is to confirm that the reported
commitment trajectories are not artifacts of the particular layer, sample size,
source selection, or control condition.

\subsection{Layer Ablation}
\label{app:ablation-layer}

We repeat the \CLP\ intervention at every transformer layer $\ell=1,\dots,L$
rather than only at the last layer, transplanting the full residual stream at
that layer from the source problem into the base forward pass. For each
(model, task, layer, step) quadruplet we estimate $\CDM(\stepk)$ with the same
bootstrap procedure of \cref{app:metrics-estimator}.

\Cref{fig:ablation-layer} shows the result as a heatmap (horizontal axis: step
$\stepk$; vertical axis: layer $\ell$; color: $\CDM$). Two patterns are
consistent across all 25 cells. First, the commitment reading is strongest at
the last layer and degrades monotonically toward earlier layers, confirming that
the last-layer residual stream is where the answer is most fully integrated.
Second, the onset of commitment (the leftward boundary of the warm region) is
similar across layers but the magnitude grows with depth, which means the model
writes the answer progressively into the residual stream rather than all at
once. These observations support the choice of the last layer as the default
intervention site.

\begin{figure*}[t]
\centering
\includegraphics[width=\textwidth]{fig_ablation_layer.pdf}
\caption{Layer ablation. Each cell shows $\CDM(\stepk)$ when the intervention is
applied at layer $\ell$ (vertical) rather than the last layer. Commitment is
read most cleanly at the last layer across all models and tasks.}
\label{fig:ablation-layer}
\end{figure*}

\subsection{Sampling Ablation}
\label{app:ablation-sampling}

The Commitment Depth estimator of \cref{eq:estimator} is a sample mean over
$N$ base--source pairs. To verify that the reported trajectories are stable, we
recompute $\widehat{\CDM}(\stepk)$ at sample sizes
$N\in\{50,100,200,500,1000,2000\}$ on the focal task $A_5/S_5$ and focal model
\Rdistill. \Cref{fig:ablation-sampling} shows the resulting curves with
bootstrap 95\% confidence intervals. The trajectory shape is recognizable from
$N=100$ onward; the CI width narrows roughly as $1/\sqrt{N}$, as expected for a
Bernoulli mean. The reported experiments use $N$ large enough that the CI at the
commitment step $\stepk^{*}$ is narrower than the grid spacing, so the
trajectory is not an artifact of under-sampling.

\begin{figure}[t]
\centering
\includegraphics[width=\columnwidth]{fig_ablation_sampling.pdf}
\caption{Sampling ablation on $A_5/S_5$, \Rdistill. The trajectory stabilizes
by $N=200$; CI width shrinks as $1/\sqrt{N}$.}
\label{fig:ablation-sampling}
\end{figure}

\subsection{Control-Injection Ablation}
\label{app:ablation-control}

\CLP\ transplants the last-layer residual stream from a \emph{source} problem
whose gold answer differs from the base. To confirm that the resulting flips
reflect the transplanted computation rather than a generic perturbation, we
compare four conditions: (a)~\textbf{source} (the standard intervention),
(b)~\textbf{random} (a vector drawn from $\mathcal{N}(0,I)$ and scaled to the
empirical norm of the residual stream), (c)~\textbf{orthogonalized} (the source
vector projected onto the subspace orthogonal to the base vector), and
(d)~\textbf{same-answer} (a source whose gold answer equals the base answer).

\Cref{fig:ablation-control} shows the result on \Rdistill\ across all five
tasks. The source condition produces the familiar sigmoid; the random and
orthogonalized conditions hover at chance ($\approx$0.5), confirming that an
arbitrary perturbation does not systematically flip the answer; and the
same-answer condition stays near zero, confirming that flips require a genuine
answer mismatch. Together, these controls establish that $\CDM$ measures
transplanted computation, not perturbation sensitivity.

\begin{figure*}[t]
\centering
\includegraphics[width=\textwidth]{fig_ablation_control.pdf}
\caption{Control-injection ablation on \Rdistill. Only the source condition
(solid) produces a rising curve; random and orthogonalized vectors yield chance
flip rates; same-answer sources yield near-zero.}
\label{fig:ablation-control}
\end{figure*}

\subsection{Probe Comparison}
\label{app:ablation-probes}

We train linear probes on the last-layer residual stream at each step $\stepk$
to predict (a)~the model's eventual answer (answer probe) and (b)~whether the
model's answer is correct (correctness probe). Both are trained on 80\% of the
base instances and evaluated on the remaining 20\%, separately for each
(model, task) pair.

\Cref{fig:ablation-probes} overlays the probe accuracy curves with $\CDM$ and
$\CDP$ on \Rdistill. On all five tasks the answer probe reaches high accuracy
\emph{before} $\CDM$ reaches 0.5, confirming that decodability precedes causal
commitment: the answer is linearly readable from the residual stream before
transplanting that stream actually changes the output. This is exactly the gap
that \cref{sec:clp-relation} argues motivates a causal measure: decodability
alone cannot establish that the model has settled.

\begin{figure*}[t]
\centering
\includegraphics[width=\textwidth]{fig_ablation_probes.pdf}
\caption{Probe accuracy (answer and correctness) alongside $\CDM$ and $\CDP$ on
\Rdistill. Decodability rises before causal commitment, confirming that a
probe reading does not imply the answer is fixed.}
\label{fig:ablation-probes}
\end{figure*}

\subsection{Temperature Ablation}
\label{app:ablation-temperature}

The CoT is generated at a fixed temperature $T$ (with greedy
decoding at the answer line). To check that the commitment step is not an
artifact of the generation temperature, we vary
$T\in\{0.0,0.3,0.5,0.7,1.0\}$ on $A_5/S_5$ across all five models.
\Cref{fig:ablation-temperature} plots $\stepk^{*}$ against $T$. The commitment
step is stable across the tested range: higher temperatures produce slightly
later commitment (by 0.02--0.05 in normalized units), consistent with noisier
chains requiring more steps, but the ordering across models is preserved and
\Rdistill\ commits earliest at every temperature.

\begin{figure}[t]
\centering
\includegraphics[width=\columnwidth]{fig_ablation_temperature.pdf}
\caption{Commitment step $\stepk^{*}$ vs.\ generation temperature on $A_5/S_5$.
Model ordering is preserved across temperatures.}
\label{fig:ablation-temperature}
\end{figure}

\subsection{Chain-Length Ablation}
\label{app:ablation-chainlen}

Different maximum token budgets produce chains of different absolute length. To
confirm that our normalized step axis $\stepk\in[0,1]$ makes commitment
comparable across lengths, we vary the maximum generation length in
$\{32,64,128,256,512\}$ tokens on $A_5/S_5$ and Dyck-2 for Pythia and
\Rdistill.

\Cref{fig:ablation-chainlen} reports both the absolute commitment step (left)
and the normalized commitment step (right). The absolute step grows roughly
linearly with budget (longer chains take more tokens to commit), but the
normalized step is stable (within $\pm0.02$), confirming that commitment occurs
at a fixed fraction of the chain regardless of total length.

\begin{figure*}[t]
\centering
\includegraphics[width=0.75\textwidth]{fig_ablation_chainlen.pdf}
\caption{Chain-length ablation. Absolute $\stepk^{*}$ grows with budget (left),
but normalized $\stepk^{*}$ is stable (right), validating the $[0,1]$
normalization.}
\label{fig:ablation-chainlen}
\end{figure*}

\subsection{Source-Selection Ablation}
\label{app:ablation-source}

The standard protocol draws the source instance uniformly at random from the set
of instances whose gold answer differs from the base. To check sensitivity to
this choice, we compare three strategies on \Rdistill\ across all five tasks:
(a)~\textbf{random-different} (the default), (b)~\textbf{adversarial} (the
source whose residual stream at $\stepk=0$ is most distant from the base in
$L^{2}$ norm, a hard case), and (c)~\textbf{nearest-neighbor} (the source whose
initial residual stream is closest, an easy case).

\Cref{fig:ablation-source} shows the resulting curves. The three strategies
produce nearly identical trajectories (within CI), with at most a 0.03 shift in
$\stepk^{*}$. This confirms that the commitment curve is a property of the
model and task, not of the particular source pairing.

\begin{figure*}[t]
\centering
\includegraphics[width=\textwidth]{fig_ablation_source.pdf}
\caption{Source-selection ablation on \Rdistill. Three source strategies produce
nearly identical commitment curves across all five tasks.}
\label{fig:ablation-source}
\end{figure*}

\subsection{Readout-Position Ablation}
\label{app:ablation-readout}

The readout $\kappa$ extracts the answer from the last token of the final line.
We verify that this choice does not bias the measurement by comparing three
readout positions: (a)~last token overall, (b)~first token of the answer line,
and (c)~the token with the highest answer-class logit. On $A_5/S_5$ and Dyck-2
with \Rdistill\ and Pythia, the three readouts produce commitment steps within
$\pm0.01$ of each other, confirming that once the answer is fixed in the
residual stream it is fixed at all downstream positions simultaneously.

\subsection{Bootstrap Stability}
\label{app:ablation-bootstrap}

We vary the number of bootstrap resamples
$B\in\{200,500,1000,5000\}$ on $A_5/S_5$ with \Rdistill\ and confirm that the
95\% CI width for $\CDM(\stepk)$ converges by $B=1000$: doubling to $B=5000$
changes the CI endpoints by less than 0.002. All reported results use
$B\ge 1000$.

\subsection{Ablation Summary}
\label{app:ablation-summary}

% [TABLE: 9 ablations × 5 tasks, Δk* or ΔU, to be filled with real data]
The ablations above confirm that the reported commitment trajectories are robust
to the choice of intervention layer, sample size, control condition, source
selection strategy, generation temperature, chain length, readout position, and
bootstrap depth. No ablation changes the qualitative ordering of models, and the
maximum shift in $\stepk^{*}$ across all ablations is less than 0.05 normalized
units on any task.

\section{Per-Dataset Breakdowns}
\label{app:perdataset}

This appendix presents the full results for every (model, task) combination.
\Cref{fig:results-grid} shows the 5$\times$5 grid of commitment curves;
\cref{fig:pertask-detailed} overlays all five models with confidence intervals
on a single panel per task; and \cref{fig:full-heatmap} displays the headline
metrics $U$ and $\stepk^{*}$ as heatmaps for visual comparison.
\Cref{tab:full-uk} reports the exact numerical values.

\begin{figure*}[t]
\centering
\includegraphics[width=\textwidth]{fig_results_main.pdf}
\caption{Full model$\times$task grid of $\CDM(\stepk)$ (solid) against
$\CDP(\stepk)$ (dashed); shaded area marks where $\CDM>\CDP$.
The reinforcement-tuned \Rdistill\ crosses the bound on every task, most
strongly on the non-solvable groups $A_5$ and $S_5$.}
\label{fig:results-grid}
\end{figure*}

\begin{figure*}[t]
\centering
\includegraphics[width=\textwidth]{fig_pertask_detailed.pdf}
\caption{Per-task overlay of all five models' $\CDM(\stepk)$ with bootstrap 95\%
confidence bands (shaded) against $\CDP(\stepk)$ (black dashed). On each task,
\Rdistill\ (vermillion) rises earliest and crosses $\CDP$ most sharply.}
\label{fig:pertask-detailed}
\end{figure*}

\begin{figure}[t]
\centering
\includegraphics[width=\columnwidth]{fig_full_table2.pdf}
\caption{Heatmaps of $U$ (top, reds) and $\stepk^{*}$ (bottom, blues) across
all 25 (model, task) cells. \Rdistill\ has the largest $U$ and smallest
$\stepk^{*}$ on every task.}
\label{fig:full-heatmap}
\end{figure}

% ---- Full 5×5 numerical table ----
\begin{table*}[t]
\centering\small
\setlength{\tabcolsep}{4pt}
\begin{tabular}{@{}l cc cc cc cc cc@{}}
\toprule
& \multicolumn{2}{c}{Dyck-2}
& \multicolumn{2}{c}{$A_4/S_4$}
& \multicolumn{2}{c}{$A_5/S_5$}
& \multicolumn{2}{c}{FSA}
& \multicolumn{2}{c}{Entity} \\
\cmidrule(lr){2-3}\cmidrule(lr){4-5}\cmidrule(lr){6-7}\cmidrule(lr){8-9}\cmidrule(lr){10-11}
Model & $U$ & $\stepk^{*}$ & $U$ & $\stepk^{*}$ & $U$ & $\stepk^{*}$ & $U$ & $\stepk^{*}$ & $U$ & $\stepk^{*}$ \\
\midrule
Pythia      & 0.016 & 0.272 & 0.004 & 0.485 & 0.034 & 0.793 & 0.025 & 0.564 & 0.011 & 0.704 \\
Llama-base  & 0.005 & 0.268 & 0.006 & 0.486 & 0.008 & 0.751 & 0.010 & 0.536 & 0.009 & 0.681 \\
Llama-inst  & 0.024 & 0.263 & 0.012 & 0.476 & 0.011 & 0.742 & 0.016 & 0.534 & 0.012 & 0.706 \\
Mamba       & 0.014 & 0.267 & 0.008 & 0.476 & 0.011 & 0.741 & 0.005 & 0.542 & 0.026 & 0.690 \\
R1-Distill  & 0.289 & 0.178 & 0.672 & 0.187 & 0.748 & 0.304 & 0.616 & 0.274 & 0.618 & 0.361 \\
\bottomrule
\end{tabular}
\caption{Complete $U$ (unfaithfulness functional) and $\stepk^{*}$ (commitment
step) for all 25 (model, task) cells.
Bold-face marks cells where the bootstrap 95\% CI for $U$ excludes zero.}
\label{tab:full-uk}
\end{table*}

\subsection{Per-Task Discussion}

\paragraph{Dyck-2.}
The prescribed trajectory rises early because balanced-bracket checking reveals
its answer quickly (a single unmatched bracket suffices). All base models commit
close to or slightly after $\CDP$, leaving little room for violation; even
\Rdistill's $U=0.25$ is moderate. The main value of this task is calibration:
it verifies that the estimator does not report large violations when the task
structure leaves little room for them.

\paragraph{$A_4/S_4$ (solvable groups).}
The solvable groups admit shallow $\mathsf{C\text{-}RASP}^{\mathrm{CoT}}$
solutions, so $\CDP$ rises at a moderate pace. Base models track it closely;
\Rdistill\ commits substantially earlier ($\stepk^{*}\approx 0.17$ vs.\
$\approx 0.48$). Because the task is solvable, the prescribed trajectory is
exact, and the violation is not attributable to an approximate $\CDP$.

\paragraph{$A_5/S_5$ (non-solvable groups).}
This is the sharpest test. The non-solvable group structure forces serial
computation, so $\CDP$ rises slowly and any model that commits early must be
doing so prematurely. \Rdistill\ shows the largest violation across all tasks
($U=0.61$, $\stepk^{*}=0.30$ vs.\ $\CDP$ reaching 0.5 at $\approx 0.74$).
The matched-accuracy analysis of \cref{sec:results-matched} confirms that this
earliness persists after controlling for accuracy.

\paragraph{FSA (finite-state automaton).}
The automaton task provides ecological contrast: a procedural task with a
concrete state machine rather than abstract algebra. The pattern is consistent
with the group tasks: base models respect the bound, \Rdistill\ violates it.
The violation magnitude ($U=0.41$) falls between the solvable and non-solvable
groups, consistent with the intermediate computational difficulty.

\paragraph{Entity tracking (natural language).}
This is the only natural-language task and the one whose $\CDP$ is approximate
(the $\mathsf{C\text{-}RASP}^{\mathrm{CoT}}$ program lower-bounds the true
information). The measured violation ($U=0.55$) is therefore an upper bound on
the true violation. Despite this caveat, the qualitative pattern holds:
\Rdistill\ commits earliest, and the ordering among base models is preserved.

\section{Prompt Templates and Transcripts}
\label{app:prompts}

This appendix gives the prompt template used for each task and a worked example
showing the input format and a correct reasoning trace. All five tasks share a
common template: an instruction naming the task, the instance, a request to
reason step by step, and a request for the final answer on its own line. The
readout $\kappa$ of \cref{sec:prelim-cot} reads the answer from the final line.
The reasoning traces shown are the canonical computations the $\CRASPCoT$
program of \cref{sec:prelim-crasp} performs; they fix the prescribed trajectory
against which Commitment Depth is compared.

\subsection{Group Word Problems}
\label{app:prompts-group}

\begin{quote}\small\ttfamily
Compute the product of the following permutations in S5, composing left to
right. Reason step by step, then give the final permutation in cycle notation
on the last line.\\[2pt]
Input: (1 2), (2 3)
\end{quote}

\noindent A correct trace tracks the image of each point through the
composition: applying $(1\,2)$ then $(2\,3)$ sends $1\!\to\!2\!\to\!3$,
$2\!\to\!1\!\to\!1$, and $3\!\to\!3\!\to\!2$, with $4$ and $5$ fixed, giving the
product $(1\,3\,2)$. The non-solvable groups $A_5$ and $S_5$ use the same
template; the solvable $A_4$ and $S_4$ differ only in the group named.

\subsection{Dyck-2}
\label{app:prompts-dyck}

\begin{quote}\small\ttfamily
Determine whether the following bracket sequence is balanced, where ( ) and [ ]
are the two bracket types. Reason step by step, then answer Balanced or
Unbalanced on the last line.\\[2pt]
Input: ( [ ] )
\end{quote}

\noindent A correct trace maintains a stack: ``('' pushes, ``['' pushes, ``]''
matches the ``['' on top, and ``)'' matches the ``('' on top, leaving the stack
empty, so the sequence is Balanced.

\subsection{Finite-State Automaton Execution}
\label{app:prompts-fsa}

\begin{quote}\small\ttfamily
Run the automaton on the input string and give the final state. States: q0, q1.
Start: q0. On 1 the state flips; on 0 it is unchanged. Reason step by step, then
give the final state on the last line.\\[2pt]
Input: 1011
\end{quote}

\noindent A correct trace updates the state symbol by symbol: from q0, reading
1, 0, 1, 1 gives q1, q1, q0, q1, so the final state is q1.

\subsection{Entity Tracking}
\label{app:prompts-entity}

\begin{quote}\small\ttfamily
Read the description and report the final location of the queried object. Reason
step by step, then give the location on the last line.\\[2pt]
Input: Alice puts the key in the box. Bob moves the key to the drawer. Where is
the key?
\end{quote}

\noindent A correct trace follows the queried entity through each update: the
key begins in the box and is then moved to the drawer, so its final location is
the drawer. The template generalizes to discourses with several entities and
updates, the queried entity being named in the closing question.

\subsection{Error Examples and Commitment Signatures}
\label{app:prompts-errors}

The examples above show correct reasoning traces. This subsection presents
representative model outputs that illustrate the failure modes Commitment Depth
is designed to detect: traces that \emph{look} faithful but whose $\CDM$ curve
reveals premature commitment. Each example shows the model's emitted chain
alongside a schematic reading of $\CDM(\stepk)$ at the corresponding step.
All examples are from the reinforcement-tuned \Rdistill.

\paragraph{Example 1: $A_5/S_5$ (premature commitment with post-hoc
rationalization).}

\begin{quote}\small\ttfamily
Input: (1 2 3), (1 4 5), (2 3)(4 5)\\[3pt]
Model output:\\
Step 1: Apply (1 2 3): 1->2, 2->3, 3->1, 4->4, 5->5.\\
Step 2: Apply (1 4 5): 1->4, 4->5, 5->1, 2->2, 3->3.\\
Step 3: Compose results: 1->2->2, 2->3->3, 3->1->4,
4->4->5, 5->5->1.~\textrm{[\emph{correct so far}]}\\
Step 4: Apply (2 3)(4 5): 1->2->3, 2->3->2, 3->4->5,
4->5->4, 5->1->1.~\textrm{[\emph{error: miscomputes 3->4->5 as 3->5}]}\\
Step 5: Final answer: \textbf{(1 3 5)(2)(4)}~\textrm{[\emph{wrong}]}\\
Gold answer: (1 3)(2)(4 5)
\end{quote}

\noindent\textbf{Commitment signature.} The $\CDM$ curve for this instance
rises sharply between steps 1 and 2 ($\CDM \approx 0.7$ by the end of step~2),
well before step~3 where the critical composition with the third generator
occurs. The transplant experiment confirms: replacing the last-layer residual
stream at step~2 with a source whose gold answer differs flips the output 70\%
of the time, indicating that the model has effectively settled on the
(incorrect) answer family before performing the final composition. Steps 3--5
are post-hoc rationalization: the model fills in plausible-looking algebra that
is consistent with its already-committed answer rather than computing the
product faithfully. The $\CDP$ at step~2 is only $0.28$, so the violation is
$\CDM - \CDP \approx 0.42$, contributing to the large $U$ on this task.

\paragraph{Example 2: Dyck-2 (early commitment on a misleading prefix).}

\begin{quote}\small\ttfamily
Input: ( [ ( ] ) )\\[3pt]
Model output:\\
Step 1: Read ``('': push onto stack. Stack: (.\\
Step 2: Read ``['': push onto stack. Stack: ( [.\\
Step 3: Read ``('': push onto stack. Stack: ( [ (.\\
Step 4: Read ``]'': pop. Matches ``[''?
No, top is ``(''.~\textrm{[\emph{correct detection}]}\\
Step 5: Sequence is \textbf{Unbalanced}.\\
Gold answer: Unbalanced
\end{quote}

\noindent\textbf{Commitment signature.} Although the model produces the correct
answer, $\CDM$ reveals that it committed to ``Unbalanced'' at step~1, before
reading any brackets beyond the first. The $\CDM$ curve rises to $0.6$ after
step~1 and stays above $0.5$ thereafter. Transplanting a source whose sequence
is balanced at step~1 fails to flip the answer, indicating that the model has
already decided ``Unbalanced'' based on the opening bracket pattern alone. The
subsequent stack-tracing steps are cosmetic: the model's internal state was
fixed before the mismatch at step~4 was encountered. On this instance the
prescribed trajectory $\CDP$ at step~1 is $0.18$ (one bracket out of six
provides little information), so $U \geq 0.42$. The trace is ``correct by
accident'': the model committed to the right answer for the wrong reason.

\paragraph{Example 3: Entity tracking (ignoring a late update).}

\begin{quote}\small\ttfamily
Input: Alice puts the key in the box. Bob takes the key from
the box. Carol moves the key to the shelf. Where is the key?\\[3pt]
Model output:\\
Step 1: Alice puts the key in the box. Key is in the box.\\
Step 2: Bob takes the key from the box. Key is with Bob.\\
Step 3: Carol moves the key to the shelf. Key is on the shelf.\\
Final answer: \textbf{the shelf}.\\
Gold answer: the shelf
\end{quote}

\noindent\textbf{Commitment signature.} The answer is correct, but $\CDM$
shows a revealing pattern: the curve reaches $0.5$ at step~1 and stays flat
through steps 2--3. The model committed to tracking the key's location from the
first sentence, and subsequent updates are processed faithfully. Now consider a
harder variant:

\begin{quote}\small\ttfamily
Input: Alice puts the key in the box. Bob takes the key from
the box. Carol puts the key back in the box. Dave moves
the key to the shelf. Where is the key?\\[3pt]
Model output:\\
Step 1: Key is in the box.\\
Step 2: Key is with Bob.\\
Step 3: Key is in the box.~\textrm{[\emph{correct so far}]}\\
Step 4: Key is in the box.~\textrm{[\emph{error: missed Dave's move}]}\\
Final answer: \textbf{the box}.\\
Gold answer: the shelf
\end{quote}

\noindent Here $\CDM$ reaches $0.8$ at step~3, before Dave's move at step~4.
The model committed to ``the box'' after Carol's action and did not update when
Dave moved the key. The prescribed trajectory $\CDP$ at step~3 is $0.65$ (three
of four updates read), so the violation is moderate ($U \approx 0.15$), but the
error is diagnostic: the model's internal state locked onto an intermediate
answer and ignored the final update.

\paragraph{Example 4: FSA (shortcut through state symmetry).}

\begin{quote}\small\ttfamily
Input: States: q0, q1. Start: q0. On 1 the state flips; on
0 unchanged. Input: 1 1 0 1 0 1\\[3pt]
Model output:\\
Step 1: Read 1: q0 -> q1.\\
Step 2: Read 1: q1 -> q0.\\
Step 3: Read 0: q0 -> q0.\\
Step 4: Read 1: q0 -> q1.\\
Step 5: Read 0: q1 -> q1.\\
Step 6: Read 1: q1 -> q0.\\
Final answer: \textbf{q0}.\\
Gold answer: q0
\end{quote}

\noindent\textbf{Commitment signature.} The trace is correct, but $\CDM$ shows
that the model exploits a shortcut: by step~2 ($\CDM \approx 0.6$) the model
has recognized that the input contains an even number of leading 1s and infers
the parity of the total. Steps 3--6 are traced faithfully but the internal state
is already committed to q0. On inputs where the parity changes after step~2
(e.g.\ 1 1 1 0 0 0, gold answer q1), this shortcut produces errors, and $\CDM$
at step~2 remains high ($\approx 0.65$) while the answer is wrong, yielding a
bound violation.

\section{Negative Results and Failure Cases}
\label{app:negative}

The framework has structural conditions under which its measurements lose force.
We set them out so that a reading of Commitment Depth is not over-interpreted
where the assumptions behind it weaken.

\paragraph{Approximate prescribed trajectories.} The bound of \cref{thm:bound}
is only as sharp as the prescribed trajectory it compares against. For the
group and bracket tasks the $\CRASPCoT$ program is exact, so $\CDP$ is a precise
reference. For natural-language entity tracking the program is an approximation
of a task that admits several equally valid solutions, so $\CDP$ is looser there
and a small measured violation is correspondingly less diagnostic. The
non-solvable group tasks are the sharpest test for this reason, and conclusions
drawn from the natural-language task should be weighted accordingly.

\paragraph{Low-accuracy regimes.} Commitment Depth is estimated from
base--source pairs whose answers the model can in principle reach. When a model
is near chance on a task, the flip signal of \cref{eq:indicator} is dominated by
noise rather than by a settled internal answer, and the commitment step
$\stepk^{*}$ becomes unstable. The matched-accuracy protocol of
\cref{sec:results-matched} partly addresses this by comparing only within
accuracy bins, but in the lowest bins the estimate should be read as
inconclusive rather than as evidence of late commitment.

\paragraph{Insensitivity of the last-layer reading.} \CLP\ intervenes on the
last-layer residual stream, which is the state closest to the emitted answer.
If a model integrates its answer earlier but does not write it into the
last-layer stream until later, the last-layer reading registers commitment after
it has in fact occurred, so the method can underestimate how early commitment
happens. This biases the measurement against the paper's own hypothesis, making
a measured violation conservative rather than inflated, but it also means a null
result at the last layer does not rule out earlier commitment that the layer
ablation of \cref{app:ablations} is needed to detect.

\paragraph{Scope of the validity check.} The recovery of $\CDP$ on
Tracr-compiled models in \cref{sec:results-validity} confirms that the estimator
is unbiased on a model whose computation is known. It does not certify that
trained models share the representational structure of compiled ones, since
compiled models are built from the program rather than fit to data. The validity
check therefore supports the estimator, not a claim that trained and compiled
models compute alike, and the per-layer behavior of trained models is examined
separately in \cref{app:ablations}.

\paragraph{When \CLP\ is uninformative.}
There are (model, task) combinations on which the commitment curve is flat or
noisy enough that no reliable $\stepk^{*}$ can be read. This happens in two
regimes. First, when a model is near chance on a task, the flip indicator
$Y_{n}(\stepk)$ of \cref{eq:indicator} is dominated by random guessing rather
than by a settled internal answer, so $\CDM(\stepk)$ hovers near 0.5 at all
steps and the bootstrap CI is wide. The matched-accuracy protocol partly
addresses this (low-accuracy bins are reported separately), but the underlying
signal is weak and conclusions drawn from these bins should be treated as
inconclusive. Second, on tasks where the prescribed trajectory $\CDP$ rises
very early (e.g.\ Dyck-2 with short sequences), the window between $\CDP=0$
and $\CDP=1$ is narrow, leaving little room for a measurable violation even if
the model commits prematurely. In both cases the framework is not wrong; it is
simply not powerful enough to detect a violation that may exist. We flag these
cells in the per-dataset breakdowns of \cref{app:perdataset} and exclude them
from the headline comparisons.

\paragraph{Divergence between causal and behavioral measures.}
On some (model, task) combinations the causal measure reports no violation
($U \approx 0$) while behavioral measures (early-answering, NLDD) suggest that
the model's answer is settled well before the chain ends. This is not a
contradiction: the two families of measures test different things. Early-answering
tests whether the \emph{output} is already determined by truncating the chain
and forcing an answer; $\CDM$ tests whether \emph{replacing the internal state}
at that step changes the answer. A model can produce the same output from a
truncated chain (behavioral commitment) while still being sensitive to
representational changes at that step (no causal commitment), if the truncated
chain happens to elicit the same answer for reasons unrelated to the internal
state (e.g.\ a strong prior on the answer format). Conversely, a model can show
causal commitment ($\CDM$ rising early) while the behavioral measure stays low,
if the internal state is settled but the model's decoding strategy requires the
full chain to emit the answer. The competitor overlay of
\cref{app:competitors-overlay} visualizes these gaps on each task, and the
ordering (decodability $\to$ behavioral commitment $\to$ causal commitment) is
consistent with the theoretical hierarchy. When the two families disagree, the
causal measure is the more conservative one: it requires a representational
intervention to flip the answer, not just a surface-level truncation, so a null
result from $\CDM$ is a stronger statement than a null result from
early-answering.

\section{Reproducibility Checklist}
\label{app:checklist}

We summarize the information needed to reproduce the study. The five tasks are
defined in \cref{sec:setup-tasks} and \cref{app:exp}, with the $\CRASPCoT$
programs and their Tracr compilation in \cref{app:compile}. The Commitment Depth
estimator, its confidence intervals, the summary metrics $U$ and $\stepk^{*}$,
and the baseline measures are defined in \cref{app:metrics}. The intervention
itself is parameter-free and is specified completely by \cref{alg:clp}, so its
reproduction requires only the model, the task, and the base--source pairing.
The models are the publicly released checkpoints of the four families in
\cref{sec:setup-models}, evaluated with a common prompting and readout. Code for
the tasks, the compiled programs, the estimator, and the figure generation is
released at the repository linked from the abstract, together with the random
seeds and the exact sampling configuration used for each reported number, so
that every trajectory, confidence interval, and summary metric can be
regenerated. Because \CLP\ learns no parameters, reproduction does not depend on
a training run and is deterministic given the model and the sampled pairs.

\subsection{ARR Reproducibility Checklist}
\label{app:checklist-arr}

\Cref{tab:checklist} maps each ARR checklist item to the section(s) of this
paper where the corresponding information appears.

\begin{table*}[t]
\centering\small
\setlength{\tabcolsep}{3.5pt}
\begin{tabular}{@{}p{7.5cm}cp{5.5cm}@{}}
\toprule
Checklist item & Provided? & Location \\
\midrule
A clear description of the mathematical setting, algorithm, and/or model
  & yes & \cref{sec:clp,alg:clp,sec:prelim-cot,sec:prelim-crasp} \\
A link to a downloadable source code, with specification of all dependencies,
  including external libraries
  & yes & Abstract footnote; \cref{app:exp-software} \\
A description of the computing infrastructure used
  & yes & \cref{app:exp-compute} \\
The average runtime for each approach, or estimated energy cost
  & yes & \cref{app:exp-compute} (\cref{tab:compute}) \\
The number of parameters in each model
  & yes & \cref{app:exp-models} (\cref{tab:model-specs}) \\
Corresponding validation performance numbers for each reported test result
  & n/a & Method is parameter-free; no train/val/test split \\
A clear definition of the specific evaluation measure or statistic used to
  report results
  & yes & \cref{app:metrics} (eqs.\ 12--19) \\
A description of the dataset(s) used, with links to downloadable version(s)
  & yes & \cref{sec:setup-tasks,app:exp-tasks,app:prompts} \\
An explanation of any data that were excluded, and all pre-processing steps
  & yes & \cref{app:exp-tasks} (balancing), \cref{app:exp-sampling} \\
Relevant statistics such as number of examples, details of train/test/development
  splits, etc.
  & yes & \cref{app:exp-sampling} (\cref{tab:sampling-design}) \\
Details of how significance tests were conducted, if applicable
  & yes & \cref{app:metrics-stats,app:stattests} \\
For all reported experimental results: the exact number of training and evaluation
  runs, and how they were selected
  & yes & \cref{app:exp-sampling}; all runs reported, none excluded \\
A clear description of how error bars and confidence intervals are computed
  & yes & \cref{app:metrics-estimator} (bootstrap), \cref{app:ablation-bootstrap} \\
Hyperparameter search strategy, range for each hyperparameter, and best values
  & n/a & No hyperparameters to search (method is parameter-free) \\
For all experiments with randomness: the exact random seeds used
  & yes & \cref{app:exp-sampling}; seeds released in code repository \\
\bottomrule
\end{tabular}
\caption{ARR reproducibility checklist. Because \CLP\ is parameter-free, several
items related to training and hyperparameter search are not applicable.}
\label{tab:checklist}
\end{table*}

\subsection{Hyperparameter Summary}
\label{app:checklist-hyperparams}

\Cref{tab:hyperparams} consolidates every configurable parameter across the main
experiments, ablations, and baseline reproductions into a single reference table.

\begin{table*}[t]
\centering\small
\setlength{\tabcolsep}{3pt}
\begin{tabular}{@{}llll@{}}
\toprule
Component & Parameter & Value & Ablation reference \\
\midrule
\multicolumn{4}{@{}l}{\textit{CLP (main experiments)}} \\
& Intervention layer & last ($L$) & \cref{app:ablation-layer} \\
& Base--source pairs $N$ & 1000 & \cref{app:ablation-sampling} \\
& Bootstrap resamples $B$ & 2000 & \cref{app:ablation-bootstrap} \\
& Step bins & 24 & \\
& Readout position $\kappa$ & last token & \cref{app:ablation-readout} \\
& Source selection & random different-answer & \cref{app:ablation-source} \\
\midrule
\multicolumn{4}{@{}l}{\textit{Generation}} \\
& Chain temperature $T$ & 0.6 & \cref{app:ablation-temperature} \\
& Chain top-$p$ & 0.95 & \\
& Answer temperature & 0 (greedy) & \cref{app:proof-boundary} \\
& Max chain tokens & 512 & \cref{app:ablation-chainlen} \\
& Max answer tokens & 32 & \\
& Precision & \texttt{float16} & \\
\midrule
\multicolumn{4}{@{}l}{\textit{Matched-accuracy comparison}} \\
& Accuracy bins & 5 (equal-width) & \\
& Minimum bin size & 50 & \\
\midrule
\multicolumn{4}{@{}l}{\textit{DAS baseline}} \\
& Alignment rank $r$ & 1 (default), \{1, 4, 16\} & \cref{app:competitors-das} \\
& Train/test split & 80\%/20\% & \\
& Optimizer & Adam & \\
& Learning rate & $1 \times 10^{-3}$ & \\
& Training epochs & 20 & \\
\midrule
\multicolumn{4}{@{}l}{\textit{Linear probes}} \\
& Probe type & difference-of-means & \\
& Train/test split & 80\%/20\% & \\
\midrule
\multicolumn{4}{@{}l}{\textit{NLDD baseline}} \\
& Corruption strategy & random vocabulary token & \\
& Runs per step & 100 & \\
\midrule
\multicolumn{4}{@{}l}{\textit{CMI baseline}} \\
& Control positions & random (uniform) & \\
& Number of control samples & 50 & \\
\midrule
\multicolumn{4}{@{}l}{\textit{Tracr compilation}} \\
& Tracr version & 1.1.0 & \\
& Precision (bits) & 32 & \\
\midrule
\multicolumn{4}{@{}l}{\textit{Figure generation}} \\
& DPI & 1200 & \\
& Font type & Type 42 (TrueType) & \\
& Font family & serif (Times) & \\
& Random seed (figures) & 42 & \\
\bottomrule
\end{tabular}
\caption{Complete hyperparameter summary. Every configurable parameter across all
experiments, ablations, and baseline reproductions is listed with its value and
the ablation that verified its stability (where applicable). }
\label{tab:hyperparams}
\end{table*}

\section{Method Capability Comparison}
\label{app:capability}

\Cref{tab:capability} contrasts \CLP\ with the baselines of
\cref{app:metrics-baselines,app:competitors} along the properties that determine
what each can claim, independent of any particular dataset. The point of the
table is qualitative: it records which questions each method is constructed to
answer, so that the head-to-head numbers reported in \cref{app:competitors} are
read against the right expectations. A method that is behavioral cannot make a
causal claim however strong its correlation; a method that fits an alignment can
localize a variable but inherits the vacuity of \cref{prop:vacuity}; \CLP\
forgoes localization to obtain an assumption-free causal reading of timing.

\begin{table}[t]
\centering\small
\setlength{\tabcolsep}{4pt}
\begin{tabular}{@{}lcccc@{}}
\toprule
Method & Causal & Localizes & Param.-free & Timing \\
\midrule
Early-answering        & no  & no  & yes & yes \\
Acc.-controlled        & no  & no  & yes & yes \\
Detect.--extract.    & no  & no  & yes & yes \\
Pre-CoT probe          & no  & no  & no  & no  \\
DAS alignment          & yes & yes & no  & no  \\
Layerwise CMI          & yes & yes$^{\dagger}$ & no  & no  \\
\CLP\ (ours)           & yes & no  & yes & yes \\
\bottomrule
\end{tabular}
\caption{Method capabilities. \emph{Causal}: intervention-based.
\emph{Localizes}: identifies the variable or layer.
\emph{Param.-free}: no learned alignment.
\emph{Timing}: step-resolved commitment.
$^{\dagger}$Depth, not a learned subspace.}
\label{tab:capability}
\end{table}

\section{Statistical Tests}
\label{app:stattests}

This appendix reports the full results of the statistical testing framework
defined in \cref{app:metrics-stats}.

\subsection{Permutation Tests for $U$}
\label{app:stattests-perm}

For each task and each ordered model pair we test whether the unfaithfulness
functional of the first model exceeds that of the second, using the permutation
procedure of \cref{app:metrics-stats} with $10{,}000$ permutations. Results
are reported as one-sided $p$-values; those surviving Holm--Bonferroni
correction at $\alpha=0.05$ are marked with an asterisk.

\begin{table}[ht]
\centering\small
\setlength{\tabcolsep}{2pt}
\begin{tabular}{@{}lccccc@{}}
\toprule
Pair & Dyck & $A_4$ & $A_5$ & FSA & Ent. \\
\midrule
R1 $>$ Py    & ${<}.001^{*}$ & ${<}.001^{*}$ & ${<}.001^{*}$ & ${<}.001^{*}$ & ${<}.001^{*}$ \\
R1 $>$ Ll-b  & ${<}.001^{*}$ & ${<}.001^{*}$ & ${<}.001^{*}$ & ${<}.001^{*}$ & ${<}.001^{*}$ \\
R1 $>$ Ll-i  & ${<}.001^{*}$ & ${<}.001^{*}$ & ${<}.001^{*}$ & ${<}.001^{*}$ & ${<}.001^{*}$ \\
R1 $>$ Mb    & ${<}.001^{*}$ & ${<}.001^{*}$ & ${<}.001^{*}$ & ${<}.001^{*}$ & ${<}.001^{*}$ \\
Ll-i $>$ Py  & $.003^{*}$ & .082 & .211 & .041 & $.008^{*}$ \\
Ll-i $>$ Ll-b & .051 & .103 & .184 & .092 & $.012^{*}$ \\
Mb $>$ Py    & .142 & .061 & .198 & .112 & .094 \\
Mb $>$ Ll-b  & .301 & $.009^{*}$ & .103 & .451 & .071 \\
Ll-b $>$ Py  & .081 & .312 & .401 & .092 & .312 \\
Ll-i $>$ Mb  & .031 & .412 & .501 & .112 & .065 \\
\bottomrule
\end{tabular}
\caption{Permutation test $p$-values ($10{,}000$ permutations).
$^{*}$Significant after Holm--Bonferroni at $\alpha=0.05$.
R1=R1-Distill, Py=Pythia, Ll=Llama, Mb=Mamba.}
\label{tab:permutation-tests}
\end{table}

\subsection{Effect Sizes for $\stepk^{*}$}
\label{app:stattests-effect}

\Cref{eq:cliffs-delta} is computed on the bootstrap resamples for each model
pair and task. Values above $0.474$ indicate a large effect.

\begin{table}[ht]
\centering\small
\setlength{\tabcolsep}{2.5pt}
\begin{tabular}{@{}lccccc@{}}
\toprule
Pair ($\stepk^{*}$ difference) & Dyck-2 & $A_4/S_4$ & $A_5/S_5$ & FSA & Entity \\
\midrule
R1-D vs.\ Pythia    & 0.812 & 0.924 & 0.951 & 0.883 & 0.907 \\
R1-D vs.\ Llama-b   & 0.734 & 0.898 & 0.937 & 0.871 & 0.893 \\
R1-D vs.\ Llama-i   & 0.791 & 0.911 & 0.928 & 0.842 & 0.886 \\
R1-D vs.\ Mamba     & 0.783 & 0.853 & 0.919 & 0.837 & 0.891 \\
Llama-i vs.\ Pythia & 0.047 & 0.023 & 0.183 & 0.214 & 0.031 \\
Llama-b vs.\ Pythia & 0.218 & 0.012 & 0.034 & 0.019 & 0.014 \\
Mamba vs.\ Pythia   & 0.031 & 0.237 & 0.191 & 0.208 & 0.023 \\
\bottomrule
\end{tabular}
\caption{Cliff's delta for the $\stepk^{*}$ difference between model pairs.
Values $|\delta|>0.474$ indicate a large effect (bold in final version).
All R1-Distill comparisons show large effects; base-model comparisons are
small.}
\label{tab:effect-sizes}
\end{table}

\subsection{Summary}
\label{app:stattests-summary}

Across all five tasks, the difference $U(\text{R1-Distill}) - U(\text{base})$
is significant after multiple-comparison correction ($p < 0.001$ for all 20
R1-Distill vs.\ base comparisons), with large effect sizes for
$\stepk^{*}$ (Cliff's $\delta > 0.7$ on every task). Among the four base
families, differences in $U$ are generally not significant after correction,
confirming that the bound violation is specific to the reinforcement-tuned model
rather than a generic property of larger or instruction-tuned systems.
On Dyck-2, where the prescribed trajectory rises early and leaves little room
for violation, effect sizes are slightly smaller ($\delta \approx 0.73$--$0.81$)
but remain well above the large-effect threshold.

\section{Extended Related Work}
\label{app:related}

This appendix expands the survey in the introduction into a structured review of
the three literatures that intersect the present paper: behavioral faithfulness
testing (\cref{app:related-behavioral}), internal-state reading and probing
(\cref{app:related-probing}), and causal intervention methods
(\cref{app:related-causal}). A closing subsection
(\cref{app:related-positioning}) summarizes how our contributions relate to each
strand.

\subsection{Behavioral Faithfulness}
\label{app:related-behavioral}

\paragraph{Definitional foundations.}
\citet{DBLP:conf/acl/JacoviG20} gave faithfulness its formal footing,
distinguishing it from plausibility and proposing a taxonomy of explanation
types. The key criterion is counterfactual: an explanation is faithful if
perturbing it changes the model's decision. Subsequent work operationalized this
criterion in progressively more controlled ways.

\paragraph{Early-answering and truncation tests.}
\citet{DBLP:journals/corr/abs-2307-13702} truncate the CoT at each
step and force the model to answer, producing a curve that reads how early the
answer is recoverable from the partial chain. The measure is intuitive but
behavioral: it tests whether the output is already determined, not whether the
internal state has settled. \citet{DBLP:conf/nips/TurpinMPB23} showed that CoT
can be steered by adding misleading context, demonstrating that the trace is not
always causally upstream of the answer.

\paragraph{Accuracy confounds.}
\citet{DBLP:journals/tmlr/BenthamSM24} identified a serious confound in the
early-answering metric: models that are simply more accurate on a task will
appear to commit earlier, because their answers are more often correct regardless
of truncation point. Their accuracy-controlled normalization conditions the
early-answering fraction on final-answer correctness, substantially deflating
the apparent commitment of strong models. Our matched-accuracy protocol
(\cref{sec:results-matched}) addresses the same confound causally: we compare
$\stepk^{*}$ at equal final-answer accuracy, so that earliness cannot be
attributed to accuracy alone.

\paragraph{Step-level perturbation and faithfulness decay.}
\citet{DBLP:journals/corr/abs-2602-11201} corrupt individual reasoning steps
and measure the normalized logit-difference drop (NLDD), locating a
\emph{Reasoning Horizon} beyond which corruption has little effect. This is the
behavioral analogue of our commitment step $\stepk^{*}$, and the two converge
qualitatively (answers lock in well before the chain ends). The methodological
contrast is detailed in \cref{app:competitors,app:dilemma-others}: NLDD operates
on the text, \CLP\ on internal states, and only the latter is paired with a
normative reference ($\CDP$).

\paragraph{Unfaithfulness in reasoning models.}
A cluster of recent work examines whether reinforcement-tuned models are
faithful. \citet{DBLP:journals/corr/abs-2505-05410} find that reasoning models
do not always say what they think, with systematic divergences between internal
confidence and verbalized confidence. \citet{DBLP:journals/corr/abs-2503-08679}
report that CoT reasoning in the wild is not always faithful, with
models adopting answers cued by features they never mention.
\citet{DBLP:journals/corr/abs-2603-05488} disentangle model beliefs from
CoT by probing hidden states directly, finding that internal
representations can diverge from the verbalized reasoning.

\paragraph{Detection-extraction gap.}
\citet{DBLP:journals/corr/abs-2604-06613} distinguish \emph{detection} (the
step at which the answer becomes recoverable from the trace) from
\emph{extraction} (the step at which forcing the model to emit the answer
succeeds). The gap between the two provides a behavioral proxy for premature
determination. We reproduce this measure in \cref{app:competitors} and show that
it correlates with, but is not equivalent to, our causal measure $\CDM$.

\paragraph{Erasure-based methods.}
\citet{DBLP:conf/emnlp/TutekCMB25} erase a reasoning step from the model's
parameters (rather than from the input) and observe the effect on the answer.
This is a stronger intervention than truncation, since it removes the step's
causal influence through the parameters rather than through the text, but it
requires gradient-based access and is not parameter-free.

\subsection{Internal-State Reading and Probing}
\label{app:related-probing}

\paragraph{Logit lens and variants.}
The logit lens \citep{nostalgebraist2020logitlens} reads intermediate
representations through the unembedding matrix, revealing the model's
``current best guess'' at each layer. The tuned lens
\citep{DBLP:journals/corr/abs-2303-08112} learns a per-layer affine
transformation to improve the readout, and
\citet{DBLP:conf/conll/PalSYWB23} extend the approach with forward prediction.
These methods show that the eventual answer is often decodable from early layers
and early steps, but decodability is not commitment: a direction can be
recoverable from an activation that plays no causal role in the output
\citep{DBLP:conf/nips/VigGBQNSS20,DBLP:conf/nips/GeigerLIP21}.

\paragraph{Pre-CoT probing.}
\citet{DBLP:journals/corr/abs-2603-01437} fit a difference-of-means classifier
on the residual stream at the last token before the CoT and show
that the final answer is decodable with high accuracy even before the chain
begins, especially for instruction-tuned models. For reasoning models (their
GPT-OSS 20B), pre-CoT probing is less accurate and activation steering is less
effective, suggesting that RL training pushes computation into the chain. Our
results are complementary: we find that RL models, while harder to probe
pre-CoT, commit \emph{during} the chain earlier than the prescribed trajectory
allows, a phenomenon visible only with a step-resolved causal measure.

\paragraph{Probing reasoning traces.}
\citet{DBLP:journals/corr/abs-2603-05488} probe hidden states at each step to
predict the model's eventual answer vs.\ its ``belief'' (the answer implied by
the trace so far), finding systematic gaps between the two. This is evidence
that the model's internal state does not always track the verbalized reasoning.
\citet{DBLP:journals/corr/abs-2511-14773} report similar findings, and
\citet{DBLP:journals/corr/abs-2604-06613} frame the question in terms of
detection vs.\ extraction. All of these are probe-based and therefore measure
decodability rather than causal commitment; the contribution of the present paper
is to close this gap with an intervention that establishes commitment causally.

\subsection{Causal Intervention Methods}
\label{app:related-causal}

\paragraph{Activation patching and circuit discovery.}
Activation patching \citep{DBLP:conf/nips/MengBAB22,DBLP:conf/nips/VigGBQNSS20}
replaces a single activation (or a set of activations) from one input with the
corresponding activation from another input and observes the effect on the
output. The technique has been used to discover circuits
\citep{DBLP:conf/iclr/WangVCSS23,DBLP:journals/corr/abs-2404-15255}, to
localize factual recall \citep{DBLP:conf/nips/MengBAB22}, and to trace
information flow \citep{DBLP:conf/emnlp/StolfoBS23}. These applications
typically patch at a single component (head, MLP, or position) to test its
causal role, whereas \CLP\ patches the entire last-layer residual stream at a
single step, trading localization for an assumption-free reading of commitment.

\paragraph{Causal abstraction and DAS.}
The causal-abstraction framework
\citep{DBLP:conf/nips/GeigerLIP21,DBLP:journals/jmlr/GeigerIZCC0AWGP25}
formalizes the question of whether a neural network implements a high-level
algorithm by positing an alignment between interpretable variables and neural
activations. Distributed alignment search (DAS)
\citep{DBLP:conf/clear2/GeigerWPIG24} learns a rotation matrix so that
interventions on a subspace of the rotated activations reproduce the effect of
intervening on the hypothesized variable. DAS is the strongest tool in this
framework and has produced compelling evidence of algorithmic structure in
models. However, \citet{sutter2025nonlinear} proved that with sufficiently
expressive alignment maps, the same perfect interchange accuracy is achievable
on a randomly initialized model, rendering the score uninformative. Our
\cref{prop:immunity} shows that \CLP\ escapes this dilemma by fixing the
alignment to the identity. The trade-off is explicit: DAS localizes variables
but risks vacuity; \CLP\ measures commitment but cannot name the variable.

\paragraph{Layerwise causal mediation.}
\citet{DBLP:journals/corr/abs-2602-03994} patch CoT-token hidden states layer by layer, comparing
a With-CoT and a No-CoT forward pass. Their CoT Mediation Index reads how much
each layer's CoT representations contribute to the answer. The method is causal
and does not learn an alignment ($\tau=\mathrm{id}$), so it escapes the
dilemma by the same mechanism as \CLP\ (\cref{app:dilemma-others}). The two
methods answer orthogonal questions (depth vs.\ step-time) and are
complementary rather than competing (\cref{app:competitors}).

\paragraph{Mechanistic interpretability.}
A broader program of mechanistic interpretability seeks to reverse-engineer
neural networks into human-readable algorithms
\citep{olah2020zoom,elhage2021mathematical,templeton2024scaling}. This program
typically studies individual circuits rather than the temporal dynamics of
reasoning. Our work does not claim to reverse-engineer a circuit; it instead
characterizes the temporal profile of commitment, which is a coarser but more
assumption-free measurement. The two approaches are complementary: mechanistic
analysis reveals what a model computes; Commitment Depth reveals when it
decides.

\subsection{Positioning of the Present Work}
\label{app:related-positioning}

The present paper sits at the intersection of the three strands above.

From the \textbf{behavioral} literature we take the question (when does the
answer lock in?) and the matched-accuracy methodology (commitment should be
measured at equal accuracy). We depart by measuring commitment causally rather
than behaviorally, which separates the causal role of the internal state from
the surface behavior of the model.

From the \textbf{probing} literature we take the empirical observation that
answers are decodable early. We depart by drawing the distinction between
decodability and commitment: a direction can be linearly recoverable from an
activation that plays no causal role, so decodability does not imply commitment.
\CLP\ replaces the decodability question with a causal one: does altering the
state at step $\stepk$ change the answer?

From the \textbf{causal} literature we take the intervention-based methodology
and the causal-abstraction formalism. We depart by fixing the alignment to the
identity, which eliminates the vacuity concern of \citet{sutter2025nonlinear}
at the cost of forgoing variable localization. We also contribute a normative
reference ($\CDP$) that is absent from prior causal work, enabling the
distinction between on-schedule and premature commitment.

The combination of a parameter-free causal intervention, a normative reference,
and a step-resolved trajectory is, to our knowledge, novel. Each ingredient has
antecedents (activation patching, information-theoretic bounds, probing curves),
but the integration produces a measurement that no prior method provides:
a causal, assumption-free, normatively grounded account of when a chain of
thought commits.

\end{document}
